\chapter{Algebra tensorowa i zewnętrzna}
\label{ch:algebra-tensorowa}
\index{tensor}

Wektor styczny i kowektor z Rozdziału~\ref{ch:rozmaitosci} są obiektami
liniowymi. Do opisu odwzorowań wieloliniowych, operatorów i
wyznaczników potrzebujemy przestrzeni, w której symbole
$v\otimes w$ spełniają dokładnie relacje wymuszone przez dwuliniowość.
Zbudujemy ją bez wyboru bazy, a dopiero potem wyprowadzimy rachunek
współrzędnych. W całym rozdziale $\mathbb K=\R$ albo $\C$;
przestrzenie mają skończony wymiar, chyba że zaznaczono inaczej.
W przypadku zespolonym wszystkie odwzorowania wieloliniowe są
$\C$-liniowe w każdym argumencie: nie stosujemy tu sprzężenia
zespolonego. W szczególności dualność algebraiczna nie jest
utożsamieniem przez iloczyn hermitowski.

\section{Iloczyn tensorowy: definicja, istnienie i jednoznaczność}
\label{sec:iloczyn-tensorowy}

\begin{definicja}[Własność uniwersalna]
\index{wlasnosc uniwersalna@Własność uniwersalna}
\label{def:tensor-universal}
\emph{Iloczynem tensorowym} przestrzeni $V,W$ nazywamy parę
$(T,\tau)$, gdzie $T$ jest przestrzenią liniową,
$\tau\colon V\times W\to T$ jest odwzorowaniem dwuliniowym,
a dla każdej przestrzeni $Z$ i każdego odwzorowania dwuliniowego
$B\colon V\times W\to Z$ istnieje dokładnie jedno odwzorowanie
liniowe $\widetilde B\colon T\to Z$ takie, że
$B=\widetilde B\circ\tau$. Piszemy $T=V\otimes W$ i
$\tau(v,w)=v\otimes w$. Elementy $v\otimes w$ nazywamy
\emph{tensorami prostymi}.
\end{definicja}

\begin{center}
\begin{tikzpicture}[>=Latex, every node/.style={font=\small}]
 \node (vw) at (0,1.8) {$V\times W$};
 \node (tensor) at (3.8,1.8) {$V\otimes W$};
 \node (z) at (3.8,0) {$Z$};
 \draw[->,thick,manifoldblue] (vw) -- node[above] {$\tau$} (tensor);
 \draw[->,thick,manifoldred] (vw) -- node[below left] {$B$} (z);
 \draw[->,thick,manifoldgreen] (tensor) -- node[right] {$\exists!\,\widetilde B$} (z);
\end{tikzpicture}
\end{center}

Samo zapisanie $v\otimes w$ nie jest definicją: najpierw trzeba
wykazać, że przestrzeń o tej własności istnieje.

\begin{uwaga}[Co znaczy ,,formalna kombinacja liniowa''?]
\label{rem:kombinacje-formalne}
Jeśli $S$ jest dowolnym zbiorem, \emph{wolna przestrzeń liniowa}
o bazie symboli $[s]$, $s\in S$, składa się z zapisów
\[
 \sum_{s\in S} c_s[s],\qquad c_s\in\mathbb K,
 \quad c_s=0\text{ poza skończoną liczbą indeksów}.
\]
To jest kombinacja \emph{formalna}: dodajemy współczynniki
stojące przy tym samym symbolu i mnożymy je przez skalary,
ale nie przypisujemy symbolom żadnych innych relacji.
Dwa takie zapisy są równe dokładnie wtedy, gdy ich
współczynniki przy każdym $[s]$ są równe. W szczególności
$[s]$ i $[s']$ są niezależnymi wektorami, jeśli $s\ne s'$.
Można myśleć o takim zapisie jako o funkcji
$s\mapsto c_s$ o \emph{skończonym nośniku}.

W dowodzie poniżej bierzemy $S=V\times W$. Jeśli
$v,v'$ są niezależne, symbole $[(v+v',w)]$, $[(v,w)]$
i $[(v',w)]$ mają różne indeksy, więc przed podzieleniem
przez relacje są niezależne. Dopiero iloraz nakazuje równość
$[(v+v',w)]+R=[(v,w)]+[(v',w)]+R$.
Równość dwóch klas $x+R=y+R$ oznacza dokładnie $x-y\in R$.
Ta ostatnia reguła wyjaśnia, dlaczego relacje dwuliniowości
pojawiają się w konstruowanej przestrzeni, a nie były
ukrytym założeniem o symbolach.
\end{uwaga}

\begin{twierdzenie}[Konstrukcja iloczynu tensorowego]
\label{thm:tensor-istnienie}
Dla dowolnych przestrzeni liniowych $V,W$ iloczyn tensorowy istnieje.
Ponadto tensory proste rozpinają $V\otimes W$.
\end{twierdzenie}
\begin{proof}
Niech $F$ będzie wolną przestrzenią liniową o bazie formalnych
symboli $[(v,w)]$, po jednym dla każdej pary $(v,w)\in V\times W$.
Jej elementy są kombinacjami z Uwagi~\ref{rem:kombinacje-formalne}.
Ta pomocnicza przestrzeń jest na ogół nieskończenie wymiarowa,
nawet gdy $V,W$ mają skończony wymiar; dopiero iloraz narzuci relacje.
Niech $R\subseteq F$ będzie podprzestrzenią rozpiętą przez
wszystkie elementy czterech rodzajów:
\begin{align*}
 &[(v+v',w)]-[(v,w)]-[(v',w)],
 &&[(av,w)]-a[(v,w)],\\
 &[(v,w+w')]-[(v,w)]-[(v,w')],
 &&[(v,aw)]-a[(v,w)],
\end{align*}
gdzie $a\in\mathbb K$. Kładziemy $T:=F/R$ oraz
$v\otimes w:=[(v,w)]+R$. Cztery relacje pokazują kolejno,
że $\tau(v,w):=v\otimes w$ jest liniowa w każdej zmiennej.
Każda klasa w $F/R$ jest skończoną kombinacją liniową klas $[(v,w)]+R$,
więc tensory proste rozpinają $T$.

Weźmy teraz dowolne odwzorowanie dwuliniowe $B\colon V\times W\to Z$.
Na $F$ istnieje dokładnie jedno odwzorowanie liniowe
$L\colon F\to Z$ o wartościach $L([(v,w)])=B(v,w)$:
określamy je na bazie i przedłużamy liniowo. Dwuliniowość $B$
oznacza, że $L$ zeruje każdy z czterech typów generatorów $R$,
a więc $R\subseteq\ker L$. Odwzorowanie $L$ przechodzi na iloraz:
$\widetilde B(x+R):=L(x)$ dla $x\in F$. Jest dobrze określone, bo
$x+R=x'+R$ implikuje $L(x-x')=0$, i spełnia
$\widetilde B(v\otimes w)=B(v,w)$. Jeśli drugie odwzorowanie liniowe
spełniałoby tę samą równość, zgadzałoby się z $\widetilde B$
na zbiorze rozpinającym $T$, a zatem wszędzie.
\end{proof}

\begin{uwaga}[Faktoryzacja przez iloraz a izomorfizm]
\label{rem:tensor-iloraz-izomorfizm}
W dowodzie wystarczyło $R\subseteq\ker L$, aby otrzymać
odwzorowanie $\widetilde B:F/R\to Z$. To jest ten sam mechanizm
przechodzenia do ilorazu, który w ogólnej postaci
porządkuje Twierdzenie~\ref{thm:modul-izomorfizm}
w Rozdziale~\ref{ch:algebra-abstrakcyjna}
(przestrzenie liniowe są modułami nad $\mathbb K$).
Pełne twierdzenie daje $F/\ker L\cong\operatorname{im}L$.
Nie utożsamiamy natomiast $F/R$ z $\operatorname{im}L$,
chyba że $R=\ker L$; do istnienia iloczynu tensorowego
taka równość nie jest potrzebna.
\end{uwaga}

\begin{twierdzenie}[Jednoznaczność]
\label{thm:tensor-jednoznacznosc}
Jeżeli $(T,\tau)$ i $(T',\tau')$ spełniają
Definicję~\ref{def:tensor-universal}, istnieje dokładnie jeden
izomorfizm spełniający
\[
 J\colon T\to T',\qquad J\circ\tau=\tau'.
\]
\end{twierdzenie}
\begin{proof}
Własność uniwersalna $(T,\tau)$ zastosowana do $\tau'$ daje
jedyne odwzorowanie liniowe $J\colon T\to T'$ z~żądaną własnością.
Podobnie $(T',\tau')$ daje $K\colon T'\to T$ z
$K\circ\tau'=\tau$. Zatem $(KJ)\circ\tau=\tau$.
Również $\id_T\circ\tau=\tau$; jednoznaczność faktoryzacji
przez $T$ wymusza $KJ=\id_T$. Tak samo $JK=\id_{T'}$.
Odwzorowanie $J$ jest więc izomorfizmem i jedynym możliwym odwzorowaniem
zachowującym wskazane obrazy par $(v,w)$.
\end{proof}

W szczególności
\begin{equation}\label{eq:tensor-biliniowosc}
 (av+bv')\otimes w=a(v\otimes w)+b(v'\otimes w),\qquad
 v\otimes(aw+bw')=a(v\otimes w)+b(v\otimes w').
\end{equation}
\emph{Każdy} element $V\otimes W$ jest skończoną sumą tensorów
prostych, ale nie musi sam być tensorem prostym.

\section{Baza, wymiar i wiele czynników}
\label{sec:baza-tensorowa}

\begin{twierdzenie}[Baza iloczynu tensorowego]
\label{thm:baza-tensorowa}
Jeżeli $e_1,\ldots,e_m$ jest bazą $V$, a
$f_1,\ldots,f_n$ bazą $W$, to
\[
 \bigl\{e_i\otimes f_j:1\leq i\leq m,\;1\leq j\leq n\bigr\}
\]
jest bazą $V\otimes W$. W szczególności
$\dim(V\otimes W)=mn$.
\end{twierdzenie}
\begin{proof}
Dla $v=\sum_i a_i e_i$ i $w=\sum_j b_j f_j$ dwuliniowość daje
\[
 v\otimes w=\sum_{i=1}^m\sum_{j=1}^n a_ib_j(e_i\otimes f_j).
\]
Ponieważ tensory proste rozpinają $V\otimes W$, wyświetlone
wektory są układem rozpinającym. Dla liniowej niezależności
wprowadźmy przestrzeń $\mathbb K^{m\times n}$ z bazą jednostek
macierzowych $E_{ij}$: jedynka na miejscu $(i,j)$ i zera poza nim. Przypisanie
$B(v,w)=(a_i b_j)_{i,j}$ jest dwuliniowe; własność uniwersalna
daje $\widetilde B\colon V\otimes W\to\mathbb K^{m\times n}$
z $\widetilde B(e_i\otimes f_j)=E_{ij}$. Jeśli
$\sum_{i,j}c_{ij}e_i\otimes f_j=0$, zastosowanie
$\widetilde B$ daje $\sum_{i,j}c_{ij}E_{ij}=0$, czyli
wszystkie $c_{ij}=0$. Tak wykazaliśmy również podany wymiar.
\end{proof}

\begin{przyklad}[Tensor prosty a suma tensorów]
Dla $V=W=\R^2$ odwzorowanie z dowodu przypisuje
$(1,2)\otimes(3,-1)$ macierz
$\begin{psmallmatrix}3&-1\\6&-2\end{psmallmatrix}$.
Każdy niezerowy tensor prosty ma tu macierz rzędu $1$:
jest ona iloczynem kolumny współrzędnych $v$ i wiersza
współrzędnych $w$. Natomiast
$e_1\otimes e_1+e_2\otimes e_2$ odpowiada macierzy
identycznościowej rzędu $2$, więc nie jest prosty.
Należy też pamiętać, że $v\otimes w$ nie oznacza pary $(v,w)$:
na przykład $(2v)\otimes w=v\otimes(2w)$.
\end{przyklad}

\begin{stwierdzenie}[Łączność i przestawianie czynników]
\label{prop:tensor-lacznosc}
Istnieją kanoniczne izomorfizmy
\[
 (V\otimes W)\otimes U\longrightarrow V\otimes(W\otimes U),
 \quad (v\otimes w)\otimes u\longmapsto v\otimes(w\otimes u),
\]
\[
 V\otimes W\longrightarrow W\otimes V,
 \quad v\otimes w\longmapsto w\otimes v.
\]
Dlatego możemy pisać $V_1\otimes\cdots\otimes V_k$ bez nawiasów.
\end{stwierdzenie}
\begin{proof}
Odwzorowanie $(v,w,u)\mapsto v\otimes(w\otimes u)$ jest
trójliniowe. Każde odwzorowanie trójliniowe na $V\times W\times U$
faktoryzuje jednoznacznie przez $(V\otimes W)\otimes U$:
najpierw przy ustalonym $u$ faktoryzujemy względem $(v,w)$;
otrzymane odwzorowania $(V\otimes W)\to Z$ zależą liniowo od $u$,
co sprawdzamy na rozpinających $v\otimes w$. Następnie
faktoryzujemy powstałe odwzorowanie dwuliniowe
$(V\otimes W)\times U\to Z$. Otrzymujemy zatem odwzorowanie w treści
stwierdzenia. Analogiczna konstrukcja w drugą stronę daje
$v\otimes(w\otimes u)\mapsto(v\otimes w)\otimes u$.
Oba złożenia są identycznością na tensorach prostych trzech
czynników, które rozpinają odpowiednie przestrzenie.
Izomorfizm przestawiający dwa czynniki pochodzi z dwuliniowego
odwzorowania $(v,w)\mapsto w\otimes v$; jego odwrotność powstaje
po ponownym przestawieniu.
Przy większej liczbie czynników dowolne dwa sposoby zmiany
nawiasowania dają to samo odwzorowanie: na każdym tensorze
prostym zachowują kolejność i wartości wszystkich czynników,
a takie tensory rozpinają dziedzinę. To uzasadnia pomijanie
nawiasów również dla czterech i więcej przestrzeni.
\end{proof}

Kładziemy $V^{\otimes0}:=\mathbb K$ oraz
$V^{\otimes k}:=V\otimes\cdots\otimes V$ dla $k\geq1$.
Jeżeli $\dim V=n$, iterowanie Twierdzenia~\ref{thm:baza-tensorowa}
daje bazę $e_{i_1}\otimes\cdots\otimes e_{i_k}$ i wymiar
\begin{equation}\label{eq:dim-potega-tensorowa}
 \dim V^{\otimes k}=n^k\qquad(k\geq0).
\end{equation}
Dla $k=0$ jest to konwencja $n^0=1$, również gdy $V=0$.

\section{Tensory mieszane i operacje na nich}
\label{sec:operacje-tensorowe}

Oznaczmy $V^*:=\operatorname{Hom}_{\mathbb K}(V,\mathbb K)$.
Jest to przestrzeń liniowych funkcjonałów, zwanych kowektorami.
Dla nieujemnych $r,s$ przyjmujemy konwencję
\begin{equation}\label{eq:typ-tensora}
 T^r_s(V):=V^{\otimes r}\otimes(V^*)^{\otimes s}.
\end{equation}
Tensory typu $(r,s)$ mają $r$ czynników wektorowych i $s$
kowektorowych; powiązane z nimi funkcje wieloliniowe
przyjmują $r$ kowektorów i $s$ wektorów. Na tensorze prostym
wyznacza je wzór
\[
 (v_1\otimes\cdots\otimes v_r\otimes
 \alpha_1\otimes\cdots\otimes\alpha_s)
 (\xi_1,\ldots,\xi_r;w_1,\ldots,w_s)
 =\prod_{a=1}^r\xi_a(v_a)\prod_{b=1}^s\alpha_b(w_b).
\]
Puste iloczyny mają wartość $1$. Dla $r=s=0$ mamy
$T^0_0(V)=\mathbb K$; odwzorowanie bez argumentów określa po prostu
skalar (dziedziną jest jednoelementowy zbiór pustych układów).

\begin{twierdzenie}[Postać wieloliniowa tensora]
\label{thm:tensor-wieloliniowy}
Dla skończenie wymiarowej przestrzeni $V$ powyższy wzór daje
izomorfizm
\[
 T^r_s(V)\cong
 \operatorname{Mult}\bigl((V^*)^r\times V^s,\mathbb K\bigr).
\]
Jeśli $e_1,\ldots,e_n$ jest bazą $V$ i $e^1,\ldots,e^n$
bazą dualną, każdy tensor ma dokładnie jedno rozwinięcie
\begin{equation}\label{eq:wspolczynniki-tensora}
 T=\sum_{\substack{i_1,\ldots,i_r\\j_1,\ldots,j_s}}
 T^{i_1\cdots i_r}{}_{j_1\cdots j_s}
 e_{i_1}\otimes\cdots\otimes e_{i_r}
 \otimes e^{j_1}\otimes\cdots\otimes e^{j_s}.
\end{equation}
\end{twierdzenie}
\begin{proof}
Rozwinięcie i jego jednoznaczność wynikają z kolejnych
zastosowań Twierdzenia~\ref{thm:baza-tensorowa}. Wzór
na działanie na argumentach jest wieloliniowy w czynnikach
tensora, więc przez własność uniwersalną definiuje odwzorowanie liniowe
$T^r_s(V)$ do wskazanej przestrzeni funkcji. Aby wykazać jego
iniektywność, podstawmy w miejsce argumentów elementy baz:
\[
 e^{i_1},\ldots,e^{i_r};\qquad e_{j_1},\ldots,e_{j_s}.
\]
Uzyskana wartość jest dokładnie współczynnikiem
$T^{i_1\cdots i_r}{}_{j_1\cdots j_s}$. Jeśli wszystkie wartości
są zerowe, znikają więc wszystkie współczynniki.
Odwrotnie, dla dowolnej funkcji wieloliniowej $F$ wstawmy te
wartości do prawej strony~\eqref{eq:wspolczynniki-tensora}.
Otrzymany tensor ma te same wartości co $F$ na wszystkich
układach elementów bazowych. Rozwinięcie argumentów i ich
wieloliniowość dają równość na dowolnych argumentach.
\end{proof}

Przykładowo $T^0_2(V)$ to formy dwuliniowe na $V$.
Związek tensorów z operatorami wymaga pilnowania,
z której przestrzeni operator \emph{przyjmuje} argument.

\begin{twierdzenie}[Tensory i przestrzenie operatorów]
\label{thm:tensor-hom}
Dla skończenie wymiarowych $V,W$ istnieje kanoniczny izomorfizm
\begin{equation}\label{eq:tensor-hom}
 \Theta_{V,W}\colon V\otimes W^*\longrightarrow
 \operatorname{Hom}_{\mathbb K}(W,V),\qquad
 \Theta_{V,W}(v\otimes\lambda)(w):=\lambda(w)v.
\end{equation}
W szczególności po zamianie nazw przestrzeni otrzymujemy
$W\otimes V^*\cong\operatorname{Hom}_{\mathbb K}(V,W)$.
\end{twierdzenie}
\begin{proof}
Dla ustalonego $w\in W$ wyrażenie $\lambda(w)v$ zależy
liniowo od $v$ oraz liniowo od $\lambda$; samo przypisanie
$w\mapsto\lambda(w)v$ jest liniowe. Dlatego odwzorowanie
$(v,\lambda)\mapsto(w\mapsto\lambda(w)v)$ ma wartości
w $\operatorname{Hom}(W,V)$ i jest dwuliniowe.
Własność uniwersalna iloczynu tensorowego wyznacza
\emph{jedno} odwzorowanie liniowe $\Theta_{V,W}$ o wzorze
\eqref{eq:tensor-hom}. Tym samym wzór działa również
na sumach: $\Theta_{V,W}(\sum_a v_a\otimes\lambda_a)(w)
=\sum_a\lambda_a(w)v_a$. Nie musimy wybierać sposobu
zapisania takiej sumy, bo $\Theta_{V,W}$ została zdefiniowana
na całej przestrzeni tensorowej.

Zbudujmy odwrotność. Wybierzmy bazę $f_1,\ldots,f_m$
przestrzeni $W$ i jej bazę dualną $f^1,\ldots,f^m$.
Dla $A\in\operatorname{Hom}(W,V)$ połóżmy
\begin{equation}\label{eq:hom-tensor-odwrotna}
 \mathcal I(A):=\sum_{j=1}^m A(f_j)\otimes f^j\in V\otimes W^*.
\end{equation}
Sprawdzimy oba złożenia. Jeśli $w\in W$, to
\[
 \Theta_{V,W}(\mathcal I(A))(w)
 =\sum_j f^j(w)A(f_j)
 =A\!\left(\sum_j f^j(w)f_j\right)=A(w),
\]
bo $w=\sum_j f^j(w)f_j$. Odwrotnie, dla prostego
tensora $v\otimes\lambda$ mamy
\[
 \mathcal I(\Theta_{V,W}(v\otimes\lambda))
 =\sum_j\lambda(f_j)v\otimes f^j
 =v\otimes\left(\sum_j\lambda(f_j)f^j\right)
 =v\otimes\lambda.
\]
Tensory proste rozpinają $V\otimes W^*$, więc drugie
złożenie też jest identycznością wszędzie. Wzór
\eqref{eq:hom-tensor-odwrotna} wygląda zależnie od bazy,
ale jest odwrotnością kanonicznej $\Theta_{V,W}$;
dlatego jego wynik od bazy nie zależy.
\end{proof}

Przy $W=V$ otrzymujemy
$T^1_1(V)=V\otimes V^*\cong\operatorname{End}(V)$.
Na tensorze prostym odpowiada temu wzór
\begin{equation}\label{eq:endo-tensor}
 v\otimes\alpha\longmapsto
 \bigl(w\longmapsto\alpha(w)v\bigr).
\end{equation}
Na przykład $e_i\otimes e^j$ przechodzi na jednostkę
macierzową $E_{ij}$: wysyła $e_j$ na $e_i$, a pozostałe
wektory bazowe na zero. W bazie $e'_i=\sum_a A^a{}_i e_a$
macierz operatora $S$ ma postać $[S]'=A^{-1}[S]A$:
zmieniają się współczynniki, nie sam operator.
Istotnie, jeśli $[v]$ i $[v]'$ oznaczają kolumny współrzędnych,
to $[v]=A[v]'$. Stąd $[Sv]'=A^{-1}[S]A[v]'$.
Baza dualna zmienia się według
$e'^i=\sum_b(A^{-1})^i{}_b e^b$, co wynika z warunku
$e'^i(e'_j)=\delta^i_j$; używamy odwrotności bez sprzężenia zespolonego.

\begin{uwaga}[Funkcja dwuliniowa a operator]
\label{rem:mult-hom}
Z Twierdzenia~\ref{thm:tensor-wieloliniowy} dla typu $(1,1)$
dostajemy
\[
 V\otimes V^*\cong\operatorname{Mult}(V^*\times V,\mathbb K).
\]
Dlaczego równie często spotyka się zapis
$\operatorname{Mult}(V\times V^*,\mathbb K)$?
Kolejność argumentów można po prostu zamienić: jeśli
$C(\lambda,v)$ jest dwuliniowa, to $B(v,\lambda):=C(\lambda,v)$
też jest dwuliniowa. Jest to odwracalna operacja, niezależna
od wyboru bazy. W obu zapisach odpowiadający operatorowi
$A\in\operatorname{Hom}(V,V)$ wzór ma postać
\begin{equation}\label{eq:mult-operator}
 B_A(v,\lambda)=\lambda(Av),
 \qquad v\in V,\quad\lambda\in V^*.
\end{equation}
Wektor $v$ jest argumentem operatora $A$; kowektor
$\lambda$ odczytuje liczbę z wyniku $Av$.

Sprawdźmy, że przy $\dim V<\infty$ wzór ten obejmuje
\emph{każdą} taką funkcję dwuliniową $B$. Wybierzmy bazę
$e_1,\ldots,e_n$ i bazę dualną $e^1,\ldots,e^n$, po czym
zdefiniujmy
\[
 A_B(v):=\sum_{i=1}^n B(v,e^i)e_i.
\]
Odwzorowanie $A_B$ jest liniowe, bo $B$ jest liniowa w pierwszym
argumencie. Dla dowolnego $\lambda\in V^*$ zachodzi
\[
 \lambda(A_Bv)=\sum_i B(v,e^i)\lambda(e_i)
 =B\!\left(v,\sum_i\lambda(e_i)e^i\right)=B(v,\lambda).
\]
Jeżeli dwa operatory dają tę samą funkcję $B$, każdy kowektor
przyjmuje na ich wartościach ten sam wynik, więc operatory
są równe. Zatem
$\operatorname{Mult}(V\times V^*,\mathbb K)
 \cong\operatorname{Hom}(V,V)$.
Na przykład tensor $u\otimes\alpha$ odpowiada operatorowi
$A(v)=\alpha(v)u$, a następnie funkcji
$B_A(v,\lambda)=\alpha(v)\lambda(u)$.
Nie utożsamiamy przy tym $V$ z $V^*$; korzystamy jedynie
z naturalnego sparowania $\lambda(v)$.
\end{uwaga}

\begin{przyklad}[Tensor prosty i nieprosty typu $(1,1)$]
Niech $V=\R^2$, niech $e_1,e_2$ będzie jego bazą, a
$e^1,e^2$ bazą dualną. Tensor
\[
 P=(e_1+e_2)\otimes(e^1-e^2)
\]
jest \emph{prosty}: ma postać jednego iloczynu tensorowego.
Odpowiada mu operator $A_P(v)=(e^1(v)-e^2(v))(e_1+e_2)$,
o macierzy
$\begin{psmallmatrix}1&-1\\1&-1\end{psmallmatrix}$.
Wszystkie jego wartości leżą na jednej prostej, więc
$\operatorname{rank}A_P=1$. Jego funkcja dwuliniowa to
$B_P(v,\lambda)=(e^1(v)-e^2(v))\lambda(e_1+e_2)$.

Z kolei tensor
$Q=e_1\otimes e^1+e_2\otimes e^2$ jest
\emph{nieprosty}. Odpowiada mu operator
$A_Q(v)=e^1(v)e_1+e^2(v)e_2=v$, czyli macierz
$\begin{psmallmatrix}1&0\\0&1\end{psmallmatrix}$ i funkcja
$B_Q(v,\lambda)=\lambda(v)$. Gdyby $Q=u\otimes\alpha$,
obraz $A_Q$ mieściłby się w prostej $\mathbb Ku$.
Jednak $A_Q(e_1)=e_1$ i $A_Q(e_2)=e_2$ są niezależne;
stąd $Q$ nie może być jednym tensorem prostym.
\end{przyklad}

\begin{uwaga}[Kolejność czynników jest istotna]
Równość $V\otimes W^*\cong\operatorname{Hom}(V,W)$
nie jest naturalnym utożsamieniem: wzór~\eqref{eq:tensor-hom}
przyjmuje argument $w\in W$ i zwraca wynik w $V$.
Do otrzymania odwzorowań z $V$ do $W$ trzeba wziąć
$W\otimes V^*$. Przestrzenie $V\otimes W^*$ oraz
$\operatorname{Hom}(V,W)$ mają wprawdzie równy wymiar,
ale równość wymiarów daje jedynie istnienie pewnego
izomorfizmu po dokonaniu dodatkowych wyborów, np.\ baz.
\end{uwaga}

\begin{twierdzenie}[Iloczyn tensorowy odwzorowań]
\label{thm:tensor-operatorow}
Dla odwzorowań liniowych $A\colon V\to V'$ i $B\colon W\to W'$
istnieje dokładnie jedno odwzorowanie liniowe
$A\otimes B\colon V\otimes W\to V'\otimes W'$ spełniające
\begin{equation}\label{eq:tensor-operatorow}
 (A\otimes B)(v\otimes w)=Av\otimes Bw.
\end{equation}
Jeżeli $A'\colon V'\to V''$ i $B'\colon W'\to W''$, to
\[
 (A'\otimes B')\circ(A\otimes B)=(A'\circ A)\otimes(B'\circ B),
 \qquad \id_V\otimes\id_W=\id_{V\otimes W}.
\]
Dla przestrzeni skończenie wymiarowych
$\operatorname{rank}(A\otimes B)=\operatorname{rank}A\,\operatorname{rank}B$.
\end{twierdzenie}
\begin{proof}
Odwzorowanie $(v,w)\mapsto Av\otimes Bw$ jest dwuliniowe:
liniowość $A$ i pierwszego miejsca $\otimes$ dają
liniowość w $v$, analogicznie w $w$. Własność uniwersalna
tworzy dokładnie jedno odwzorowanie o wzorze~\eqref{eq:tensor-operatorow}.
Obie strony równości złożeń przyłożone do $v\otimes w$
dają $A'Av\otimes B'Bw$. Tensory proste rozpinają dziedzinę,
więc odwzorowania są równe. Identyczność sprawdzamy tak samo.

Obraz $A\otimes B$ jest rozpięty przez wszystkie $Av\otimes Bw$,
więc równa się $(\operatorname{im}A)\otimes(\operatorname{im}B)$,
rozumianemu jako podprzestrzeń $V'\otimes W'$.
Aby uzasadnić to ostatnie włączenie bez ukrytego założenia,
dopełnijmy bazy obu obrazów do baz $V'$ i $W'$;
Twierdzenie~\ref{thm:baza-tensorowa} pokazuje, że iloczyny
wektorów z mniejszych baz pozostają liniowo niezależne
w większym iloczynie tensorowym. Ich liczba to iloczyn rang.
\end{proof}

Jeśli macierze $A,B$ mają współczynniki $a_{pi},b_{qj}$,
to w bazach $e_i\otimes f_j$ współczynnik macierzy
$A\otimes B$ o indeksach $(p,q),(i,j)$ jest $a_{pi}b_{qj}$.
Jest to \emph{iloczyn Kroneckera}. Pary porządkujemy
leksykograficznie: $(1,1),(1,2),\ldots,(2,1),\ldots$.
W takim porządku na przykład
\[
 A=\begin{pmatrix}1&2\\0&1\end{pmatrix},\quad
 B=\begin{pmatrix}2&0\\0&3\end{pmatrix},\qquad
 A\otimes B=
 \begin{pmatrix}
 2&0&4&0\\0&3&0&6\\0&0&2&0\\0&0&0&3
 \end{pmatrix}.
\]

Iloczyn tensorów mieszanych wymaga ustalenia kolejności.
Jeśli $S\in T^r_s(V)$ i $T\in T^{r'}_{s'}(V)$, ich surowy
iloczyn ma kolejno bloki $V^{\otimes r}$, $(V^*)^{\otimes s}$,
$V^{\otimes r'}$, $(V^*)^{\otimes s'}$. Przenosimy trzeci blok
przed drugi kanonicznym przestawieniem czynników i tak otrzymany
element $T^{r+r'}_{s+s'}(V)$ oznaczamy także $S\otimes T$.
Wewnątrz bloków zachowujemy kolejność. Nie pojawia się znak:
jest to zwykły iloczyn tensorowy, a nie iloczyn zewnętrzny.

\begin{definicja}[Kontrakcja]
\index{kontrakcja@Kontrakcja}
\label{def:kontrakcja}
\emph{Kontrakcja} usuwa jedną parę: czynnik $V$ i czynnik $V^*$,
łącząc je przez ewaluację $\alpha(v)$. W najprostszym przypadku
jest to odwzorowanie $C\colon V\otimes V^*\to\mathbb K$,
$C(v\otimes\alpha)=\alpha(v)$. W~ogólności dostajemy
$C_{a,b}\colon T^r_s(V)\to T^{r-1}_{s-1}(V)$, gdy $r,s\geq1$.
Indeksy spełniają $1\leq a\leq r$ i $1\leq b\leq s$.
Na tensorze prostym mnożymy przez $\alpha_b(v_a)$,
usuwamy czynniki $v_a$ i $\alpha_b$, a pozostałe zachowujemy
w pierwotnej kolejności. Przykładowo kontrakcja pierwszej pary ma współczynniki
\[
 (C_{1,1}T)^{i_2\cdots i_r}{}_{j_2\cdots j_s}
 =\sum_{a=1}^{\dim V}T^{a i_2\cdots i_r}{}_{a j_2\cdots j_s}.
\]
\end{definicja}

\begin{stwierdzenie}[Kontrakcja i ślad]
\label{prop:kontrakcja-slad}
Kontrakcja jest liniowa i nie zależy od wyboru bazy. Dla
$S=\sum_{i,j}S^i{}_j e_i\otimes e^j\in T^1_1(V)$ zachodzi
\[
 C(S)=\sum_i S^i{}_i=\operatorname{tr}S,
\]
gdzie ostatni ślad dotyczy operatora z~\eqref{eq:endo-tensor}.
\end{stwierdzenie}
\begin{proof}
Ewaluacja $(v,\alpha)\mapsto\alpha(v)$ jest dwuliniowa,
więc ma jednoznaczne przedłużenie do $V\otimes V^*$.
Pozostałe czynniki przestawiamy izomorfizmami ze
Stwierdzenia~\ref{prop:tensor-lacznosc} i stosujemy tę ewaluację
w wybranej parze. Ponieważ definicja używa jedynie wartości
kowektora na wektorze, wynik nie zależy od bazy.
Dla współrzędnych $C(e_i\otimes e^j)=e^j(e_i)=\delta_i^j$,
więc $C(S)=\sum_i S^i{}_i$. Z~\eqref{eq:endo-tensor}
współczynniki $S^i{}_j$ są współczynnikami macierzy operatora,
a powyższa suma jest jej śladem.
\end{proof}

\begin{przyklad}[Iloczyn i kontrakcja]
Niech $S=e_1\otimes e^1+2e_2\otimes e^2$ na $\R^2$.
Wtedy $S(e_1)=e_1$, $S(e_2)=2e_2$, a $C(S)=3$.
Dla $u\otimes v\otimes\alpha\in T^2_1(V)$ kontrakcja
$v$ z $\alpha$ daje $\alpha(v)u$, zaś kontrakcja
$u$ z $\alpha$ daje $\alpha(u)v$. Wybór pary ma znaczenie.
\end{przyklad}

\section{Algebra tensorowa i symetryzacja}
\label{sec:algebra-tensorowa}

Wyjaśnijmy użyte dalej słowa. \emph{Algebra} to przestrzeń
liniowa z dwuliniowym mnożeniem; u nas mnożenie będzie
łączne i będzie miało element jednostkowy $1$.
\emph{Suma prosta} $\bigoplus_{k\geq0}U_k$ zawiera ciągi
$(u_0,u_1,\ldots)$ z tylko skończenie wieloma niezerowymi
składnikami. \emph{Ideał obustronny} $I$ algebry $A$ to
podprzestrzeń taka, że $ai,ia\in I$ dla wszystkich
$a\in A$, $i\in I$. Dzięki temu na klasach w $A/I$
można określić mnożenie $(a+I)(b+I)=ab+I$:
zamiana $a$ na $a+i$ i $b$ na $b+j$ dodaje do
iloczynu wyrazy $aj+ib+ij\in I$. Algebra jest
\emph{stopniowana}, gdy $A=\bigoplus_{k\geq0}A^k$
i $A^kA^\ell\subseteq A^{k+\ell}$. Element należący do $A^k$ nazywamy
jednorodnym stopnia $k$. Homomorfizm algebr z jedynką to
odwzorowanie liniowe zachowujące mnożenie i jedynkę.

\begin{samepage}
\begin{definicja}[Algebra tensorowa]
\index{algebra tensorowa@Algebra tensorowa}
\label{def:algebra-tensorowa}
\emph{Algebra tensorowa} przestrzeni $V$ to suma prosta
\[
 T(V):=\bigoplus_{k=0}^{\infty}V^{\otimes k},
 \qquad V^{\otimes0}=\mathbb K.
\]
Elementy tej sumy mają tylko skończenie wiele niezerowych
składników. Mnożenie tensorów jednorodnych definiujemy przez
\emph{konkatenację}:
\[
 (v_1\otimes\cdots\otimes v_k)
 \cdot(w_1\otimes\cdots\otimes w_\ell)
 :=v_1\otimes\cdots\otimes v_k\otimes w_1\otimes\cdots\otimes w_\ell.
\]
Rozszerzamy je dwuliniowo na $T(V)$; jedynką jest $1\in\mathbb K$.
W stopniu $0$ mnożenie jest zwykłym mnożeniem przez skalar.
Definicja na tensorach prostych przechodzi na całe potęgi tensorowe
dzięki iterowanej własności uniwersalnej: jest liniowa w każdym
czynniku. Skończoność nośnika sum zapewnia, że wynik znów należy do $T(V)$.
\end{definicja}
\end{samepage}

\begin{twierdzenie}[Łączność i swoboda algebry tensorowej]
\label{thm:algebra-tensorowa-universal}
$T(V)$ jest łączną algebrą z jedynką, na ogół nieprzemienną.
Dla każdej łącznej algebry $\mathcal A$ z jedynką i każdego
odwzorowania liniowego $f\colon V\to\mathcal A$ istnieje dokładnie jeden
homomorfizm algebr z jedynką $\widehat f\colon T(V)\to\mathcal A$
spełniający $\widehat f|_V=f$.
\end{twierdzenie}
\begin{proof}
Na tensorach prostych obu stron równości
$(ab)c=a(bc)$ otrzymujemy ten sam tensor przez zapisanie
wszystkich czynników po kolei. Stwierdzenie~\ref{prop:tensor-lacznosc}
pozwala pominąć nawiasy; dwuliniowość rozszerza równość
na sumy skończone. Podobnie $1\cdot a=a\cdot1=a$.

Dla ustalonego $k\geq1$ odwzorowanie
$(v_1,\ldots,v_k)\mapsto f(v_1)\cdots f(v_k)$
jest $k$-liniowe, ponieważ mnożenie w $\mathcal A$ jest
dwuliniowe. Iterowana własność uniwersalna daje jedyne odwzorowanie
liniowe $f_k\colon V^{\otimes k}\to\mathcal A$ spełniające
$f_k(v_1\otimes\cdots\otimes v_k)=f(v_1)\cdots f(v_k)$.
W stopniu $0$ kładziemy $f_0(a)=a1_{\mathcal A}$.
Definiujemy $\widehat f$ na całej sumie prostej jako sumę
$f_k$ po jej skończenie wielu niezerowych składowych.
Dla tensorów prostych zachodzi
$\widehat f(a\cdot b)=\widehat f(a)\widehat f(b)$;
dwuliniowość dowodzi tego dla dowolnych $a,b$.
Jedynka też jest zachowana. Jeśli $g$ jest drugim takim
homomorfizmem, to na $v_1\otimes\cdots\otimes v_k$
musi przyjmować wartość $g(v_1)\cdots g(v_k)
=f(v_1)\cdots f(v_k)$; te tensory i $1$ rozpinają $T(V)$.
Stąd $g=\widehat f$.
\end{proof}

Dla niezależnych $e_1,e_2$ tensory $e_1\otimes e_2$ i
$e_2\otimes e_1$ są różnymi elementami $V^{\otimes2}$.
Wymuszenie przemienności prowadzi do algebry symetrycznej,
a wymuszenie $v\wedge v=0$ do algebry zewnętrznej.

\begin{definicja}[Symetryzator i antysymetryzator]
\index{symetryzator i antysymetryzator@Symetryzator i antysymetryzator}
\label{def:sym-alt}
Dla $k\geq1$ permutacja $\sigma\in S_k$ działa na
$V^{\otimes k}$ przez
\[
 P_\sigma(v_1\otimes\cdots\otimes v_k)
 :=v_{\sigma^{-1}(1)}\otimes\cdots\otimes v_{\sigma^{-1}(k)}.
\]
Odwzorowania te istnieją na mocy iterowanej własności uniwersalnej.
Definiujemy projekcje
\begin{equation}\label{eq:sym-alt}
 \operatorname{Sym}_k:=\frac1{k!}\sum_{\sigma\in S_k}P_\sigma,
 \qquad
 \operatorname{Alt}_k:=\frac1{k!}\sum_{\sigma\in S_k}
\operatorname{sgn}(\sigma)P_\sigma.
\end{equation}
W stopniu $0$ przyjmujemy
$\operatorname{Sym}_0=\operatorname{Alt}_0=\id_{\mathbb K}$.
\end{definicja}

\begin{twierdzenie}[Własności obu projekcji]
\label{thm:sym-alt}
Dla $t\in V^{\otimes k}$ zachodzi
\[
 \operatorname{Sym}_k^2=\operatorname{Sym}_k,\qquad
 \operatorname{Alt}_k^2=\operatorname{Alt}_k,
\]
a ich obrazy są dokładnie przestrzeniami tensorów spełniających,
odpowiednio, $P_\sigma t=t$ albo
$P_\sigma t=\operatorname{sgn}(\sigma)t$ dla każdego $\sigma$.
\end{twierdzenie}
\begin{proof}
Dzięki wyborowi $\sigma^{-1}$ w definicji
$P_\sigma P_\tau=P_{\sigma\tau}$; sprawdzamy to bezpośrednio
na tensorach prostych. Dla ustalonej $\tau$ podstawienie
$\rho=\tau\sigma$ w sumach~\eqref{eq:sym-alt} daje
\[
 P_\tau\operatorname{Sym}_k=\operatorname{Sym}_k,
 \qquad
 P_\tau\operatorname{Alt}_k
 =\operatorname{sgn}(\tau)\operatorname{Alt}_k.
\]
Tak samo, zmieniając indeks po prawej stronie, otrzymujemy
$\operatorname{Sym}_kP_\tau=\operatorname{Sym}_k$ oraz
$\operatorname{Alt}_kP_\tau=\operatorname{sgn}(\tau)\operatorname{Alt}_k$.
Każdy element pierwszego obrazu jest więc niezmienniczy,
a drugiego zmienia znak zgodnie z parzystością permutacji.
Jeżeli odwrotnie $P_\sigma t=t$ dla wszystkich $\sigma$,
to $\operatorname{Sym}_k t=\frac1{k!}\sum_\sigma t=t$.
Jeśli $P_\sigma t=\operatorname{sgn}(\sigma)t$, to
$\operatorname{Alt}_k t=
\frac1{k!}\sum_\sigma\operatorname{sgn}(\sigma)^2t=t$.
Zastosowanie projekcji do jej obrazu nic więc nie zmienia,
co daje obie równości idempotencji.
\end{proof}

Na przykład
\[
 \operatorname{Sym}_2(u\otimes v)
 =\tfrac12(u\otimes v+v\otimes u),\qquad
 \operatorname{Alt}_2(u\otimes v)
 =\tfrac12(u\otimes v-v\otimes u).
\]
Dla $u=v$ drugi tensor znika. W stopniu $k=1$ obie
projekcje są identycznością.

\begin{przyklad}[Rozkład tensora dwukowariantnego]
Niech $e^1,e^2$ będzie bazą $V^*$ i
$t=e^1\otimes e^2+3e^2\otimes e^1$.
Wówczas $t(e_1,e_2)=1$, $t(e_2,e_1)=3$, a
\[
 \operatorname{Sym}_2(t)
 =2(e^1\otimes e^2+e^2\otimes e^1),\qquad
 \operatorname{Alt}_2(t)
 =-e^1\otimes e^2+e^2\otimes e^1.
\]
Ich suma jest równa $t$. Pierwszy składnik jest symetryczny,
drugi zmienia znak po zamianie argumentów.
\end{przyklad}

Równość $t=\operatorname{Sym}_2t+\operatorname{Alt}_2t$
zachodzi dla każdego tensora drugiego stopnia, bo w sumie
$\tfrac12(\id+P_{(12)})+\tfrac12(\id-P_{(12)})$ zostaje $\id$.
Nie jest to ogólny rozkład na dwie części dla $k\geq3$.
Na przykład dla trzech niezależnych wektorów bazowych
współczynnik $e_1\otimes e_2\otimes e_3$ w
$(\operatorname{Sym}_3+\operatorname{Alt}_3)
(e_1\otimes e_2\otimes e_3)$ wynosi $2/6$, a nie $1$.

\begin{definicja}[Algebra symetryczna]
\index{algebra symetryczna@Algebra symetryczna}
\label{def:algebra-symetryczna}
Niech $J$ będzie obustronnym ideałem w $T(V)$ generowanym
przez wszystkie $u\otimes v-v\otimes u$. Iloraz
$S(V):=T(V)/J$ jest \emph{algebrą symetryczną}.
Jego składową stopnia $k$ oznaczamy $S^k(V)$.
Jest ona równa $V^{\otimes k}/(J\cap V^{\otimes k})$.
Istotnie, $J$ jest sumą swoich części jednorodnych:
w każdym elemencie $a(u\otimes v-v\otimes u)b$ rozwijamy
$a,b$ na skończone sumy składników jednorodnych. Każda
otrzymana część nadal należy do $J$. To uzasadnia stopniowanie ilorazu.
\end{definicja}

\begin{stwierdzenie}[Baza algebry symetrycznej]
\label{prop:sym-baza}
Dla bazy $e_1,\ldots,e_n$ i $k\geq0$ klasy
$e_1^{a_1}\cdots e_n^{a_n}$, gdzie
$a_i\geq0$ i $a_1+\cdots+a_n=k$, tworzą bazę $S^k(V)$.
Zatem
\[
 \dim S^k(V)=\binom{n+k-1}{k}\quad(n\geq1),
 \qquad S^k(V)\cong\operatorname{im}\operatorname{Sym}_k\quad(k\geq1).
\]
\end{stwierdzenie}
\begin{proof}
Niech $q\colon T(V)\to S(V)$ będzie ilorazem.
Ponieważ $q(e_i)q(e_j)=q(e_j)q(e_i)$, każdą klasę tensora
prostego można uporządkować do postaci
$e_1^{a_1}\cdots e_n^{a_n}$; te jednomiany rozpinają.
Zdefiniujmy homomorfizm $\pi\colon S(V)\to
\mathbb K[x_1,\ldots,x_n]$ przez $\pi(q(e_i))=x_i$.
Istnieje, bo odwzorowanie z algebry tensorowej wysyłające $e_i$
na $x_i$ zeruje każdy generator $J$. Z drugiej strony
przypisanie $x_i\mapsto q(e_i)$ wyznacza homomorfizm
$\rho\colon\mathbb K[x_1,\ldots,x_n]\to S(V)$:
na wielomianie podstawiamy za zmienne te przemienne elementy.
Oba złożenia ustalają wszystkie generatory i jedynkę,
więc są identycznościami. Niezależność i rozpinanie
wynikają teraz z niezależności zwykłych jednomianów.
Liczbę rozwiązań $a_1+\cdots+a_n=k$ otrzymujemy przez
rozmieszczenie $n-1$ przegród wśród $n+k-1$ pozycji,
czyli $\binom{n+k-1}{k}$.

Na koniec w stopniu $k$ zachodzi $q_kP_\sigma=q_k$,
zatem $q_k\operatorname{Sym}_k=q_k$.
Symetryzator zeruje generatory jednorodnej części ideału
$J$: zamiana dwóch sąsiednich czynników nie zmienia
średniej po wszystkich permutacjach. Przechodzi zatem
na odwzorowanie $S^k(V)\to\operatorname{im}\operatorname{Sym}_k$.
W jedną stronę jego złożenie z $q_k$ daje symetryzator,
w drugą $q_k\operatorname{Sym}_k=q_k$; są to izomorfizmy odwrotne.
Dokładniej: ograniczamy $q_k$ do obrazu symetryzatora,
na którym $\operatorname{Sym}_k$ jest identycznością.
W stopniu $0$ obie przestrzenie to $\mathbb K$ i izomorfizm
jest identycznością. Jeśli $V=0$, to $S^0(V)=\mathbb K$
i $S^k(V)=0$ dla $k>0$; wzoru z symbolem Newtona dla $n\ge1$
nie stosujemy w tym przypadku.
\end{proof}

Dla $V=\R^2$ mamy np. bazę
$S^2(V)=\langle e_1^2,e_1e_2,e_2^2\rangle$,
podczas gdy $\dim V^{\otimes2}=4$. W~algebrze symetrycznej
$e_1e_2=e_2e_1$, ale w $T(V)$ odpowiednie tensory są różne.

\section{Algebra zewnętrzna i własność uniwersalna}
\label{sec:algebra-zewnetrzna}

\begin{definicja}[Algebra zewnętrzna]
\index{algebra zewnetrzna@Algebra zewnętrzna}
\label{def:algebra-zewnetrzna}
Niech $I$ będzie obustronnym ideałem algebry $T(V)$
generowanym przez $v\otimes v$ dla wszystkich $v\in V$.
\emph{Algebra zewnętrzna} jest ilorazem
\[
 \Lambda(V):=T(V)/I,
 \qquad \Lambda^k(V):=V^{\otimes k}/(I\cap V^{\otimes k}).
\]
Iloczyn w ilorazie zapisujemy $\wedge$, a obraz
$v_1\otimes\cdots\otimes v_k$ oznaczamy
$v_1\wedge\cdots\wedge v_k$.
\end{definicja}

Ideał $I$ jest jednorodny: jego generatory mają stopień $2$,
a każdy iloczyn z tensorami jednorodnymi ma określony stopień.
Stąd podana definicja $\Lambda^k(V)$ jest zgodna z rozkładem
$\Lambda(V)=\bigoplus_{k\geq0}\Lambda^k(V)$.
W stopniach $0$ i $1$ ideał nie ma elementów niezerowych, dlatego
$\Lambda^0(V)=\mathbb K$ i $\Lambda^1(V)=V$.

\begin{stwierdzenie}[Relacje i znaki]
\label{prop:wedge-znak}
W $\Lambda(V)$ zachodzą
\begin{equation}\label{eq:wedge-relacje}
 v\wedge v=0,\qquad u\wedge v=-v\wedge u.
\end{equation}
Jeżeli $\alpha\in\Lambda^k(V)$ i $\beta\in\Lambda^\ell(V)$,
to
\begin{equation}\label{eq:wedge-graded}
 \alpha\wedge\beta=(-1)^{k\ell}\beta\wedge\alpha.
\end{equation}
Mnożenie $\wedge$ jest dwuliniowe, łączne i ma jedynkę $1$.
\end{stwierdzenie}
\begin{proof}
Pierwsza równość~\eqref{eq:wedge-relacje} jest relacją
z definicji ideału. Po podstawieniu $u+v$ dostajemy
$0=(u+v)\wedge(u+v)=u\wedge v+v\wedge u$;
stąd druga równość. Iloraz algebry łącznej z jedynką przez
ideał obustronny dziedziczy łączność i dwuliniowość.
Ponieważ ideał składa się z elementów stopni co najmniej $2$,
nie zawiera jedynki i jej klasa pozostaje jedynką.

Dla prostych $\alpha=v_1\wedge\cdots\wedge v_k$ i
$\beta=w_1\wedge\cdots\wedge w_\ell$ przejście od
$\alpha\wedge\beta$ do $\beta\wedge\alpha$ wymaga
$k\ell$ zamian miejsc sąsiednich wektorów. Każda zamiana
zmienia znak, co dowodzi~\eqref{eq:wedge-graded} dla
prostych wielowektorów. Te rozpinają składowe jednorodne,
a dwuliniowość kończy dowód.
\end{proof}

\begin{twierdzenie}[Własność uniwersalna algebry zewnętrznej]
\label{thm:wedge-algebra-universal}
Niech $\mathcal A$ będzie łączną algebrą z jedynką,
a $f\colon V\to\mathcal A$ odwzorowaniem liniowym takim, że
$f(v)^2=0$ dla każdego $v\in V$. Istnieje dokładnie jeden
homomorfizm algebr z jedynką
$\bar f\colon\Lambda(V)\to\mathcal A$ spełniający
$\bar f(v)=f(v)$ dla $v\in\Lambda^1(V)=V$.
\end{twierdzenie}
\begin{proof}
Z Twierdzenia~\ref{thm:algebra-tensorowa-universal} mamy
jedyny homomorfizm $\widehat f\colon T(V)\to\mathcal A$
przedłużający $f$. Dla każdego generatora ideału $I$
zachodzi $\widehat f(v\otimes v)=f(v)^2=0$.
Jądro homomorfizmu jest ideałem obustronnym, więc zawiera
cały $I$. Odwzorowanie
$\bar f([t]):=\widehat f(t)$ na $T(V)/I$ jest zatem
dobrze określone i jest homomorfizmem. Każdy homomorfizm
przedłużający $f$ musi na klasie
$v_1\wedge\cdots\wedge v_k$ przyjmować wartość
$f(v_1)\cdots f(v_k)$; te klasy wraz z jedynką rozpinają
$\Lambda(V)$. Stąd jednoznaczność.
\end{proof}

Przypomnijmy związek alternowania ze znakiem. Jeśli odwzorowanie
wieloliniowe $b$ znika przy powtórzeniu argumentów, to rozwinięcie
$0=b(\ldots,u+v,\ldots,u+v,\ldots)$ w dwóch wskazanych miejscach
daje
\[
 b(\ldots,u,\ldots,v,\ldots)
 =-b(\ldots,v,\ldots,u,\ldots).
\]
Wyrazy z dwoma $u$ i dwoma $v$ znikają. Każda permutacja jest
złożeniem transpozycji, więc przestawienie argumentów mnoży wynik
przez jej znak. Odwrotnie, zmiana znaku przy każdej transpozycji
daje $b=-b$ przy równych argumentach, czyli $b=0$ nad $\R$ lub $\C$.
Stąd w tym rozdziale „alternujący” i „antysymetryczny” są równoważne
dla odwzorowań wieloliniowych.

\begin{twierdzenie}[Własność uniwersalna stopnia $k$]
\label{thm:wedge-multilinear}
Dla $k\geq1$ odwzorowanie
\[
 (v_1,\ldots,v_k)\longmapsto v_1\wedge\cdots\wedge v_k
\]
jest $k$-liniowe i alternujące: znika, gdy dwa jego argumenty
są równe. Dla każdej przestrzeni $Z$ i każdego alternującego
$k$-liniowego odwzorowania $b\colon V^k\to Z$ istnieje dokładnie jedno
odwzorowanie liniowe $\bar b\colon\Lambda^k(V)\to Z$ takie, że
\[
 \bar b(v_1\wedge\cdots\wedge v_k)=b(v_1,\ldots,v_k).
\]
\end{twierdzenie}
\begin{proof}
Wieloliniowość wynika z definicji mnożenia i z liniowego
włączenia $V=\Lambda^1(V)$. Jeśli dwa wektory są równe,
przesuwamy jeden przez sąsiednie czynniki, korzystając
z~\eqref{eq:wedge-relacje}, aż znajdą się obok siebie.
Wtedy ich iloczyn jest zerowy.

Przez iterowaną własność iloczynu tensorowego odwzorowanie $b$
wyznacza jedyne odwzorowanie liniowe
$\widetilde b\colon V^{\otimes k}\to Z$ z
$\widetilde b(v_1\otimes\cdots\otimes v_k)=b(v_1,\ldots,v_k)$.
Jednorodna część $I\cap V^{\otimes k}$ jest rozpięta
przez elementy postaci
$x\otimes v\otimes v\otimes y$, gdzie $x,y$ są tensorami
prostymi, a łączny stopień wynosi $k$. Na każdym z nich
$\widetilde b$ znika, bo w $b$ występują dwa jednakowe
argumenty. Wobec tego $\widetilde b$ przechodzi na iloraz
$\Lambda^k(V)$ i określa $\bar b$.
Wielowektory proste rozpinają ten iloraz, więc faktoryzacja
jest jednoznaczna.
\end{proof}

\begin{twierdzenie}[Antysymetryzacja a iloczyn zewnętrzny]
\label{thm:wedge-alt}
Niech $q_k\colon V^{\otimes k}\to\Lambda^k(V)$ będzie odwzorowaniem
ilorazowym. Istnieje kanoniczny izomorfizm
\[
 j_k\colon\Lambda^k(V)\longrightarrow
 \operatorname{im}\operatorname{Alt}_k\subseteq V^{\otimes k},
 \qquad j_k(v_1\wedge\cdots\wedge v_k)
 =\operatorname{Alt}_k(v_1\otimes\cdots\otimes v_k).
\]
W szczególności $j_2(u\wedge v)=
\tfrac12(u\otimes v-v\otimes u)$. Przy tym utożsamieniu,
dla tensorów antysymetrycznych $a$ stopnia $k$ i $b$ stopnia
$\ell$ ich iloczyn zewnętrzny odpowiada
$\operatorname{Alt}_{k+\ell}(a\otimes b)$.
\end{twierdzenie}
\begin{proof}
Jeżeli tensor prosty zawiera kolejno $v\otimes v$, to
transpozycja tych dwóch miejsc pozostawia go bez zmian,
ale zmienia znak jego antysymetryzacji zgodnie z
Twierdzeniem~\ref{thm:sym-alt}. Ponieważ pracujemy nad
$\R$ lub $\C$, jego antysymetryzacja jest zerowa.
To samo dotyczy wielokrotności i sum takich tensorów,
zatem $\operatorname{Alt}_k$ zeruje $I\cap V^{\otimes k}$
i przechodzi na wskazane odwzorowanie $j_k$.

W ilorazie permutacja czynników zmienia znak o
$\operatorname{sgn}(\sigma)$, na mocy~\eqref{eq:wedge-relacje}.
Dla wielowektora prostego mamy więc
\[
 q_kj_k(v_1\wedge\cdots\wedge v_k)
 =\frac1{k!}\sum_{\sigma\in S_k}
 \operatorname{sgn}(\sigma)^2
 (v_1\wedge\cdots\wedge v_k)
 =v_1\wedge\cdots\wedge v_k.
\]
Wynika stąd iniektywność $j_k$. Z drugiej strony
$j_kq_k=\operatorname{Alt}_k$ na tensorach prostych,
a więc wszędzie, co dowodzi surjektywności na obraz.
Jeśli $a=j_k(\alpha)$ i $b=j_\ell(\beta)$, to
$j_{k+\ell}(\alpha\wedge\beta)
=\operatorname{Alt}_{k+\ell}(a\otimes b)$,
bo $q_k(a)=\alpha$, $q_\ell(b)=\beta$ i stosujemy właśnie
sprawdzoną równość $j_{k+\ell}q_{k+\ell}=\operatorname{Alt}_{k+\ell}$.
\end{proof}

\begin{uwaga}[Skąd bierze się czynnik $1/k!$?]
\label{rem:normalizacja-wedge}
Obraz $u\wedge v$ w \emph{algebrze zewnętrznej} to klasa
$u\otimes v$. Dopiero jego kanoniczny reprezentant w
\emph{przestrzeni tensorów antysymetrycznych}, dany przez $j_2$,
ma postać $\tfrac12(u\otimes v-v\otimes u)$. W szczególności
nie wolno utożsamiać tych dwóch zapisów bez podania
normalizacji. Dla kowektorów standardowe odwzorowanie
$\Lambda^k(V^*)$ na funkcje alternujące na $V^k$
przyjmiemy z normalizacją przez wyznacznik w
Sekcji~\ref{sec:wyznacznik}.
\end{uwaga}

\section{Baza i wymiar potęg zewnętrznych}
\label{sec:baza-zewnetrzna}

\begin{twierdzenie}[Baza potęgi zewnętrznej]
\label{thm:wedge-baza}
Niech $e_1,\ldots,e_n$ będzie bazą $V$. Dla $0\leq k\leq n$
rodzina
\[
 e_{i_1}\wedge\cdots\wedge e_{i_k},
 \qquad 1\leq i_1<\cdots<i_k\leq n,
\]
stanowi bazę $\Lambda^k(V)$; dla $k=0$ jest to wektor $1$.
Dla $k>n$ zachodzi $\Lambda^k(V)=0$. W szczególności
\begin{equation}\label{eq:wedge-wymiar}
 \dim\Lambda^k(V)=\binom nk,\qquad
 \dim\Lambda(V)=\sum_{k=0}^n\binom nk=2^n.
\end{equation}
\end{twierdzenie}
\begin{proof}
Rozwijamy każdy czynnik $v_a$ w bazie $e_i$ i korzystamy
z wieloliniowości iloczynu zewnętrznego. Otrzymujemy sumę
wielowektorów $e_{j_1}\wedge\cdots\wedge e_{j_k}$.
Jeśli indeks się powtarza, składnik znika, bo dwa jednakowe
wektory można ustawić obok siebie. Pozostałe porządkujemy
rosnąco, a każda zamiana sąsiadów daje minus. Rodzina
z tezy rozpięła więc $\Lambda^k(V)$; dla $k>n$ wszystkie
ciągi $k$ indeksów mają powtórzenie, więc
cała przestrzeń jest zerowa.

Aby udowodnić niezależność, nie wystarczy samo porządkowanie.
Dla rosnącego indeksu $I=(i_1<\cdots<i_k)$ i
$v_b=\sum_j a_{jb}e_j$ definiujemy
\begin{equation}\label{eq:wedge-wspolrzedne-minor}
 D_I(v_1,\ldots,v_k)
 :=\sum_{\sigma\in S_k}\operatorname{sgn}(\sigma)
 \prod_{b=1}^k a_{i_{\sigma(b)},b}.
\end{equation}
W każdym składniku dla ustalonego $b$ występuje dokładnie
jeden współczynnik wektora $v_b$, więc $D_I$ jest
wieloliniowa. Zamiana dwóch argumentów przestawia dwie
kolumny: permutacje w sumie parują się ze znakami przeciwnymi.
Dlatego $D_I$ jest alternująca. Własność uniwersalna z
Twierdzenia~\ref{thm:wedge-multilinear} daje funkcjonał
$\overline D_I\colon\Lambda^k(V)\to\mathbb K$.
Na uporządkowanym wielowektorze bazowym $e_{j_1}\wedge\cdots
\wedge e_{j_k}$ wartość $\overline D_I$ wynosi $1$,
gdy $I=(j_1,\ldots,j_k)$, oraz $0$ w przeciwnym razie.
Jeżeli zatem suma $\sum_J c_J e_{j_1}\wedge\cdots\wedge e_{j_k}$
jest zerowa, zastosowanie kolejno $\overline D_I$ daje
$c_I=0$ dla każdego $I$. Rodzina jest liniowo niezależna.
Indeksów rosnących długości $k$ jest $\binom nk$;
sumę wymiarów obliczamy z dwumianu Newtona.
\end{proof}

\begin{przyklad}[Współczynniki $u\wedge v$]
Dla $V=\R^3$, $u=e_1+2e_2$ i $v=e_2+e_3$ mamy
\[
 u\wedge v=(e_1+2e_2)\wedge(e_2+e_3)
 =e_1\wedge e_2+e_1\wedge e_3+2e_2\wedge e_3.
\]
Współczynnik przy $e_i\wedge e_j$ jest wyznacznikiem
$\begin{psmallmatrix}u^i&v^i\\u^j&v^j\end{psmallmatrix}
=u^iv^j-u^jv^i$, zgodnie z~\eqref{eq:wedge-wspolrzedne-minor}.
\end{przyklad}

\begin{przyklad}[Nie każdy dwuwektor jest prosty]
W $\Lambda^2(\R^4)$ weźmy
$\omega=e_1\wedge e_2+e_3\wedge e_4$. Wówczas
\[
 \omega\wedge\omega
 =2e_1\wedge e_2\wedge e_3\wedge e_4\ne0.
\]
Natomiast $(u\wedge v)\wedge(u\wedge v)=0$ dla każdego
$u,v$, gdyż $u$ występuje dwa razy. Zatem $\omega$ nie
jest pojedynczym iloczynem $u\wedge v$. Pokazuje to, że
również w algebrze zewnętrznej potrzebujemy sum wielowektorów
prostych.
\end{przyklad}

\section{Potęgi zewnętrzne i wyznacznik}
\label{sec:wyznacznik}

\begin{twierdzenie}[Potęga zewnętrzna odwzorowania]
\label{thm:potega-zewn-operatora}
Dla liniowego $A\colon V\to W$ i $k\geq1$ istnieje dokładnie jedno odwzorowanie
liniowe
\[
 \Lambda^k A\colon\Lambda^k(V)\to\Lambda^k(W),\qquad
 (\Lambda^k A)(v_1\wedge\cdots\wedge v_k)
 =Av_1\wedge\cdots\wedge Av_k.
\]
Dla $B\colon W\to U$ zachodzi
$\Lambda^k(B\circ A)=\Lambda^k B\circ\Lambda^k A$
i $\Lambda^k\id_V=\id_{\Lambda^k(V)}$.
Przyjmujemy również $\Lambda^0A=\id_{\mathbb K}$; wtedy powyższe
prawa obowiązują także dla $k=0$, a pusty iloczyn zewnętrzny oznacza $1$.
\end{twierdzenie}
\begin{proof}
Odwzorowanie $(v_1,\ldots,v_k)\mapsto
Av_1\wedge\cdots\wedge Av_k$ jest wieloliniowe,
bo $A$ jest liniowe, a iloczyn zewnętrzny wieloliniowy.
Gdy dwa argumenty są równe, odpowiadające im obrazy też
są równe, więc iloczyn znika. Twierdzenie~\ref{thm:wedge-multilinear}
daje jedyne odwzorowanie $\Lambda^k A$.
Na wielowektorze prostym obie strony wzoru na złożenie
dają $BA v_1\wedge\cdots\wedge BA v_k$.
Wielowektory proste rozpinają $\Lambda^k(V)$, co dowodzi
równości na całej przestrzeni. Argument dla identyczności
jest taki sam.
\end{proof}

\begin{definicja}[Wyznacznik operatora]
\index{wyznacznik operatora@Wyznacznik operatora}
\label{def:det-wedge}
Jeżeli $\dim V=n$, to $\Lambda^n(V)$ jest jednowymiarowa
na mocy Twierdzenia~\ref{thm:wedge-baza}. Dla operatora
$A\colon V\to V$ jego \emph{wyznacznik} jest jedynym skalarem
$\det A\in\mathbb K$, dla którego
\begin{equation}\label{eq:det-top-wedge}
 \Lambda^n A=(\det A)\id_{\Lambda^n(V)}.
\end{equation}
Dla $V=0$ przyjmujemy $\det\id_V=1$; odpowiada temu
wyznacznik pustej macierzy równy $1$.
\end{definicja}

\begin{twierdzenie}[Wzór Leibniza i własności wyznacznika]
\label{thm:det-wlasnosci}
Jeśli $Ae_j=\sum_{i=1}^n a_{ij}e_i$, to
\begin{equation}\label{eq:det-leibniz}
 \det A=\sum_{\sigma\in S_n}\operatorname{sgn}(\sigma)
 \prod_{j=1}^n a_{\sigma(j),j}.
\end{equation}
Ponadto
\[
 \det\id_V=1,\qquad \det(B\circ A)=\det B\,\det A,
 \qquad \det A\ne0\ \Longleftrightarrow\ A\text{ jest odwracalne}.
\]
Wartość $\det A$ nie zależy od bazy; wyznacznik jest
liniowy w każdej kolumnie macierzy i zmienia znak po
zamianie dwóch kolumn.
\end{twierdzenie}
\begin{proof}
Niech $\varepsilon=e_1\wedge\cdots\wedge e_n\ne0$.
Z~\eqref{eq:det-top-wedge} i wieloliniowości mamy
\[
 (\det A)\varepsilon
 =Ae_1\wedge\cdots\wedge Ae_n
 =\sum_{i_1,\ldots,i_n}
   a_{i_1,1}\cdots a_{i_n,n}
   e_{i_1}\wedge\cdots\wedge e_{i_n}.
\]
Wyrazy z powtarzającym się indeksem znikają. Pozostałe
ciągi $(i_1,\ldots,i_n)$ są permutacjami
$(\sigma(1),\ldots,\sigma(n))$, a przestawienie odpowiadającego
iloczynu zewnętrznego do $\varepsilon$ daje znak
$\operatorname{sgn}(\sigma)$. Porównanie współczynników
przy niezerowym $\varepsilon$ daje~\eqref{eq:det-leibniz}.

Twierdzenie~\ref{thm:potega-zewn-operatora} w stopniu $n$
daje $\Lambda^n(B\circ A)=\Lambda^n B\circ\Lambda^n A$.
Ponieważ każde z tych odwzorowań jest mnożeniem przez odpowiedni
skalar, otrzymujemy $\det(BA)=\det B\det A$;
tak samo $\det\id_V=1$. Jeśli $A$ ma odwrotność,
$1=\det(A^{-1}A)=\det(A^{-1})\det A$, więc $\det A\ne0$.
Jeżeli $A$ nie ma odwrotności, kolumny $Ae_1,\ldots,Ae_n$
są liniowo zależne. Jedna z nich jest kombinacją pozostałych,
a po rozwinięciu odpowiadającego czynnika każdy wyraz
$Ae_1\wedge\cdots\wedge Ae_n$ zawiera powtórzony wektor
i znika. Stąd $\det A=0$.

Definicja~\eqref{eq:det-top-wedge} mówi o skalarze działania
operatora na samej jednowymiarowej przestrzeni
$\Lambda^n(V)$, dlatego nie zależy od wyboru jej bazy.
Wreszcie $Ae_1\wedge\cdots\wedge Ae_n$ jest liniowy
w każdym czynniku i antysymetryczny przy ich zamianie;
jego współczynnik przy $\varepsilon$, czyli wyznacznik,
ma takie same własności względem kolumn.
\end{proof}

\begin{przyklad}[Pole zorientowane w $\R^2$]
Niech $u=(3,1)$ i $v=(1,2)$ w standardowej bazie.
Wówczas
\[
 u\wedge v=(3e_1+e_2)\wedge(e_1+2e_2)
 =(3\cdot2-1\cdot1)e_1\wedge e_2
 =5e_1\wedge e_2.
\]
Wartość $5$ jest polem równoległoboku ze znakiem określonym
przez kolejność $(u,v)$; pole geometryczne wynosi $|5|$.
Zmiana kolejności wektorów daje $v\wedge u=-5e_1\wedge e_2$.
\end{przyklad}

\begin{figure}[H]
\centering
\begin{tikzpicture}[>=Latex,scale=1.15]
 \coordinate (O) at (0,0);
 \coordinate (U) at (3,1);
 \coordinate (V) at (1,2);
 \coordinate (UV) at (4,3);
 \fill[manifoldblue!13] (O)--(U)--(UV)--(V)--cycle;
 \draw[->,gray!65] (-0.3,0)--(4.65,0) node[right] {$x$};
 \draw[->,gray!65] (0,-0.3)--(0,3.6) node[above] {$y$};
 \draw[thick,manifoldblue!75] (U)--(UV)--(V);
 \draw[->,very thick,manifoldred] (O)--(U)
   node[pos=0.82,below] {$u=(3,1)$};
 \draw[->,very thick,manifoldgreen] (O)--(V)
   node[pos=0.72,left] {$v=(1,2)$};
 \fill[manifoldblue] (UV) circle (1.6pt) node[above right] {$u+v$};
 \node[fill=white,inner sep=2pt] at (2.1,1.45)
   {$\det(u,v)=5$};
\end{tikzpicture}
\caption{Iloczyn zewnętrzny w $\R^2$ koduje pole i orientację
równoległoboku. Zmiana kolejności wektorów zmienia znak.}
\label{fig:wedge-pole}
\end{figure}

\begin{twierdzenie}[Minory jako macierz $\Lambda^k A$]
\label{thm:minory-lambda}
Niech $A\colon\mathbb K^n\to\mathbb K^m$ ma macierz
$(a_{ij})$. W bazach uporządkowanych z
Twierdzenia~\ref{thm:wedge-baza} współczynnik macierzy
$\Lambda^k A$ w wierszu $I=(i_1<\cdots<i_k)$ i kolumnie
$J=(j_1<\cdots<j_k)$ jest wyznacznikiem podmacierzy
$A_{I,J}=(a_{i_a,j_b})_{a,b=1}^k$.
\end{twierdzenie}
\begin{proof}
Rozwijamy
$Ae_{j_1}\wedge\cdots\wedge Ae_{j_k}$ według
$Ae_{j_b}=\sum_i a_{i,j_b}e_i$. Przy bazowym
$e_{i_1}\wedge\cdots\wedge e_{i_k}$ pojawiają się tylko
składniki, których $k$ indeksów stanowi permutację $I$.
Po ustawieniu ich w kolejności rosnącej dostajemy
\[
 \sum_{\sigma\in S_k}\operatorname{sgn}(\sigma)
 \prod_{b=1}^k a_{i_{\sigma(b)},j_b}
 =\det A_{I,J}
\]
na mocy wzoru~\eqref{eq:det-leibniz} dla macierzy $k\times k$.
Nie występują żadne inne składniki przy wskazanym
wektorze bazowym, więc to jest szukany współczynnik.
\end{proof}

\begin{wniosek}[Wzór Cauchy'ego--Bineta]
\label{cor:cauchy-binet}
Dla macierzy $A$ o wymiarach $m\times n$ i $B$ o wymiarach
$n\times p$, oraz rosnących indeksów $I\subseteq\{1,\ldots,m\}$,
$J\subseteq\{1,\ldots,p\}$ długości $k$, gdzie $0\leq k\leq\min(m,p)$,
zachodzi
\begin{equation}\label{eq:cauchy-binet}
 \det\bigl((AB)_{I,J}\bigr)
 =\sum_{\substack{K\subseteq\{1,\ldots,n\}\\ |K|=k}}
 \det A_{I,K}\,\det B_{K,J}.
\end{equation}
\end{wniosek}
\begin{proof}
Z Twierdzenia~\ref{thm:potega-zewn-operatora} mamy
$\Lambda^k(AB)=(\Lambda^k A)(\Lambda^k B)$.
Twierdzenie~\ref{thm:minory-lambda} mówi, że macierze
trzech występujących tu operatorów mają za współczynniki
odpowiednie minory. Współczynnik iloczynu macierzy w
wierszu $I$ i kolumnie $J$ to suma po wszystkich
$k$-elementowych indeksach pośrednich $K$; jest ona
dokładnie prawą stroną~\eqref{eq:cauchy-binet}.
Dla $k=0$ obie strony są równe $1$. Gdy $k>n$,
nie ma indeksów $K$ i suma jest zerowa, zgodnie z tym,
że $\Lambda^k(\mathbb K^n)=0$.
\end{proof}

Oznaczamy przez $\operatorname{Alt}^k(V;\mathbb K)$ przestrzeń
alternujących $k$-liniowych funkcji $V^k\to\mathbb K$.
Dla $k=0$ jest to $\mathbb K$.

\begin{twierdzenie}[Kowektory zewnętrzne jako odwzorowania alternujące]
\label{thm:dual-zewnetrzny}
Dla skończenie wymiarowego $V$ istnieje kanoniczny izomorfizm
\[
 \mathcal E_k\colon\Lambda^k(V^*)
 \longrightarrow\operatorname{Alt}^k(V;\mathbb K),
\]
ustalony na iloczynach prostych przez
\begin{equation}\label{eq:wedge-form-det}
 \mathcal E_k(\alpha_1\wedge\cdots\wedge\alpha_k)
 (v_1,\ldots,v_k)
 =\det\bigl(\alpha_a(v_b)\bigr)_{a,b=1}^k.
\end{equation}
W szczególności $\Lambda^k(V^*)\cong(\Lambda^k V)^*$.
Dla kowektorów $\alpha,\beta$ otrzymujemy
\begin{equation}\label{eq:wedge-kowektory}
 \mathcal E_2(\alpha\wedge\beta)(u,v)
 =\alpha(u)\beta(v)-\alpha(v)\beta(u).
\end{equation}
Dla $k=0$ przyjmujemy $\mathcal E_0=\id_{\mathbb K}$ i wyznacznik
pustej macierzy równy $1$. Dla $k>\dim V$ obie strony izomorfizmu
są przestrzeniami zerowymi.
\end{twierdzenie}
\begin{proof}
Wyznacznik w~\eqref{eq:wedge-form-det} jest wieloliniowy
i alternujący w wierszach, a zatem w $\alpha_a$;
Twierdzenie~\ref{thm:wedge-multilinear}, zastosowane do $V^*$,
zapewnia dobrze określone odwzorowanie liniowe $\mathcal E_k$.
Jest ono alternujące w argumentach $v_b$, bo zamiana dwóch
kolumn wyznacznika zmienia znak, i wieloliniowe w tych
argumentach. Jeśli $e^1,\ldots,e^n$ jest bazą dualną,
to wartość $\mathcal E_k(e^{i_1}\wedge\cdots\wedge e^{i_k})$
na $(e_{j_1},\ldots,e_{j_k})$, dla indeksów rosnących,
wynosi $\delta_{I,J}$. Wektory bazowe w
$\Lambda^k(V^*)$ przechodzą więc na bazę dualną do
wektorów bazowych $\Lambda^k V$: daje to izomorfizm
$\Lambda^k(V^*)\cong(\Lambda^k V)^*$.
Każde odwzorowanie alternujące $V^k\to\mathbb K$ wyznacza
funkcjonał na $\Lambda^k V$ przez
Twierdzenie~\ref{thm:wedge-multilinear}, więc powyższy
izomorfizm obejmuje dokładnie wszystkie takie odwzorowania.
Przy $k=2$ wzór wyznacznika $2\times2$ daje
równość~\eqref{eq:wedge-kowektory}.
\end{proof}

\begin{uwaga}[Dwie zgodne normalizacje]
Utożsamienie $j_2$ z Twierdzenia~\ref{thm:wedge-alt} wysyła
$\alpha\wedge\beta$ na
$\tfrac12(\alpha\otimes\beta-\beta\otimes\alpha)$.
Natomiast $\mathcal E_2$ wysyła go na \emph{dwukrotność}
tego antysymetrycznego odwzorowania dwuliniowego: wzór
\eqref{eq:wedge-kowektory} nie zawiera czynnika $1/2$.
Ogólnie, po odczytaniu tensorów kowariantnych jako funkcji wieloliniowych,
zachodzi $\mathcal E_k=k!\,j_k$: suma w wyznaczniku ma te same
składniki ze znakami co antysymetryzacja, lecz bez dzielenia przez $k!$.
Po utożsamieniu elementów $\Lambda^k(V^*)$ z
odwzorowaniami alternującymi za pomocą $\mathcal E_k$, iloczyn
$\omega\wedge\eta$ stopni $k$ i $\ell$ jest dany przez
\begin{equation}\label{eq:wedge-alt-forms}
 (\omega\wedge\eta)(v_1,\ldots,v_{k+\ell})
 =\frac1{k!\ell!}\sum_{\sigma\in S_{k+\ell}}
 \operatorname{sgn}(\sigma)
 \omega(v_{\sigma(1)},\ldots,v_{\sigma(k)})
 \eta(v_{\sigma(k+1)},\ldots,v_{\sigma(k+\ell)}).
\end{equation}
Dla iloczynów kowektorów prostych jest to rozwinięcie
wyznacznika względem pierwszych $k$ wierszy: każdemu
wyborowi $k$ kolumn odpowiada $k!\ell!$ permutacji
wewnątrz dwóch bloków. Ponieważ takie iloczyny rozpinają
przestrzenie zewnętrzne, wzór zachodzi zawsze.
Równoważnie współczynnik przed
$\operatorname{Alt}_{k+\ell}(\omega\otimes\eta)$ jest równy
$(k+\ell)!/(k!\ell!)$, gdy $\omega,\eta$ są już funkcjami
alternującymi. Wynika to bezpośrednio z
$\mathcal E_m=m!j_m$ i wzoru na mnożenie dla $j_m$.
To wyjaśnia, dlaczego wzór w modelu tensorów antysymetrycznych
i wzór dla form mają różne współczynniki, mimo że opisują tę samą algebrę.
\end{uwaga}

\section{Tensory w punkcie rozmaitości}
\label{sec:tensory-rozmaitosc}

Dla gładkiej rozmaitości $M$ i $p\in M$ bierzemy teraz
$V=T_pM$ jako rzeczywistą przestrzeń liniową, a
$V^*=T_p^*M$ z Sekcji~\ref{sec:przestrzen-styczna}
i~\ref{sec:przestrzen-kostyczna}. Tensor typu $(r,s)$
\emph{w punkcie} $p$ jest elementem $T^r_s(T_pM)$;
to pojedynczy obiekt w jednym włóknie.

\begin{stwierdzenie}[Iloczyn tensorowy wiązek]
\label{prop:tensor-wiazek}
Niech $E,F\to M$ będą wiązkami wektorowymi rzędów
odpowiednio $a,b$. Przestrzenie $(E\otimes F)_p:=E_p\otimes F_p$
tworzą wiązkę wektorową rzędu $ab$ nad $M$. Jeśli
$g_{ji}$ i $h_{ji}$ są macierzami przejścia wiązek
$E$ i $F$, to macierzami przejścia $E\otimes F$
są $g_{ji}\otimes h_{ji}$.
\end{stwierdzenie}

\begin{proof}
Na obszarze, gdzie $E$ i $F$ mają bazy lokalne
$e_1,\ldots,e_a$ oraz $f_1,\ldots,f_b$, ich $ab$
iloczynów $e_u\otimes f_v$ stanowi bazę każdego włókna
$E_p\otimes F_p$ na mocy twierdzenia o bazie iloczynu
tensorowego. Gdy zmienimy obie bazy, wzór na współczynniki
iloczynu prostego pokazuje, że nowe współczynniki otrzymujemy
przez przemnożenie starych macierzą $g_{ji}\otimes h_{ji}$;
jej elementy są iloczynami gładkich funkcji. Na przecięciu
trzech obszarów własność złożenia z
Twierdzenia~\ref{thm:tensor-operatorow} daje
\[
 (g_{kj}\otimes h_{kj})(g_{ji}\otimes h_{ji})
   =(g_{kj}g_{ji})\otimes(h_{kj}h_{ji})
   =g_{ki}\otimes h_{ki}.
\]
Jest to warunek sklejania; macierze są odwracalne,
bo ich odwrotnościami są $g_{ji}^{-1}\otimes h_{ji}^{-1}$.
Konstrukcja wiązki z takich danych, podana w
Sekcji~\ref{sec:wiazki-wektorowe}, kończy dowód.
\end{proof}

\begin{definicja}[Pole tensorowe]
\index{pole tensorowe@Pole tensorowe}
\label{def:pole-tensorowe}
\emph{Pole tensorowe} typu $(r,s)$ przypisuje każdemu
$p\in M$ tensor $T_p\in T^r_s(T_pM)$ w sposób gładki.
Dokładnie: jest gładkim przekrojem wiązki
\begin{equation}\label{eq:wiazka-tensorowa}
 T^r_s M:=(TM)^{\otimes r}\otimes(T^*M)^{\otimes s}
 \longrightarrow M.
\end{equation}
Przy wykładniku $0$ rozumiemy czynnik jako wiązkę
trywialną $M\times\R$.
\end{definicja}

W lokalnej mapie $x^1,\ldots,x^n$ bazami włókien są
$\partial_i:=\partial/\partial x^i$ oraz dualne
$\dd x^j$, określone przez $\dd x^j(\partial_i)=\delta_i^j$.
Pole ma zatem postać
\begin{equation}\label{eq:pole-tensorowe-lokalnie}
 T=\sum_{\substack{i_1,\ldots,i_r\\j_1,\ldots,j_s}}
 T^{i_1\cdots i_r}{}_{j_1\cdots j_s}(x)
 \partial_{i_1}\otimes\cdots\otimes\partial_{i_r}
 \otimes\dd x^{j_1}\otimes\cdots\otimes\dd x^{j_s}.
\end{equation}
Warunek gładkości oznacza dokładnie, że wszystkie funkcje
$T^{i_1\cdots i_r}{}_{j_1\cdots j_s}(x)$ są gładkie.
W innej mapie bazowe wektory i kowektory się zmieniają;
przepisanie tej samej wartości $T_p$ w nowych bazach
narzuca współczynnikom prawa transformacji, które teraz rozpiszemy.
Niech $x$ i $y$ będą dwiema mapami wokół $p$, a
\[
 A^a{}_i=\frac{\partial y^a}{\partial x^i}(p),\qquad
 B^j{}_b=\frac{\partial x^j}{\partial y^b}(p),\qquad B=A^{-1}.
\]
Reguła łańcucha i dualność dają
\[
 \partial_{x^i}=\sum_a A^a{}_i\partial_{y^a},\qquad
 \dd x^j=\sum_b B^j{}_b\dd y^b.
\]
Podstawiając te równości w każdy czynnik
\eqref{eq:pole-tensorowe-lokalnie}, otrzymujemy
\begin{equation}\label{eq:tensor-zmiana-wspolrzednych}
 (T_y)^{a_1\cdots a_r}{}_{b_1\cdots b_s}
 =\sum_{I,J}
 \left(\prod_{\mu=1}^r A^{a_\mu}{}_{i_\mu}\right)
 \left(\prod_{\nu=1}^s B^{j_\nu}{}_{b_\nu}\right)
 (T_x)^{i_1\cdots i_r}{}_{j_1\cdots j_s}.
\end{equation}
Wszystkie wartości odnoszą się do tego samego punktu $p$;
zmienia się jedynie jego opis i baza. Każdy górny indeks wnosi
macierz Jacobiego $A$, a każdy dolny jej odwrotność $B$.
Odwracając zamianę map, odzyskujemy stare współczynniki.
Ponieważ $A,B$ są gładkie, gładkość współczynników w jednej
mapie pociąga ją w każdej innej.

Wiązka~\eqref{eq:wiazka-tensorowa} powstaje
przez wielokrotne zastosowanie
Stwierdzenia~\ref{prop:tensor-wiazek} do $TM$ i $T^*M$.

Typ $(0,0)$ to zwykła funkcja gładka, typ $(1,0)$ to
pole wektorowe z Rozdziału~\ref{ch:pola}, a typ $(0,2)$
to gładko zmieniająca się rodzina form dwuliniowych
na przestrzeniach stycznych. Iloczyn tensorowy i kontrakcja
działają \emph{punkt po punkcie}. Pozostawiają pola gładkie,
bo w mapach ich współczynniki są sumami iloczynów
gładkich funkcji współczynnikowych.

\begin{przyklad}[Tensor identycznościowy]
Operator $\id_{T_pM}$ z~\eqref{eq:endo-tensor} odpowiada
w dowolnej bazie $e_i$ tensorowi
\[
 \mathsf I_p=\sum_i e_i\otimes e^i\in T^1_1(T_pM).
\]
Jego zapis w mapie to $\sum_i\partial_i\otimes\dd x^i$.
Nie zależy od mapy: po przyłożeniu do dowolnego $v$
odzyskujemy $\sum_i e^i(v)e_i=v$, a izomorfizm
$T^1_1(T_pM)\cong\operatorname{End}(T_pM)$ jest
iniektywny. Tensor $\mathsf I$ jest zatem globalnym
gładkim polem tensorowym, mimo że zapisaliśmy go lokalnie.
\end{przyklad}

\begin{przyklad}[Tensor jako iloczyn skalarny]
Na $\R^2$ pole
$g=\dd x\otimes\dd x+\dd y\otimes\dd y$
przypisuje każdemu punktowi ten sam iloczyn skalarny.
Dla $u=a\partial_x+b\partial_y$ i
$v=c\partial_x+d\partial_y$ mamy
$g(u,v)=ac+bd$. Oznacza to, że tensor typu $(0,2)$
można czytać jako przepis przyjmujący \emph{dwa}
wektory styczne w jednym punkcie i zwracający liczbę.
Sama lista jego współczynników zmieni się przy zamianie
współrzędnych, ale wartość $g(u,v)$ pozostanie ta sama.
Na jednowymiarowym przykładzie widać konkretny czynnik:
dla $x>0$ i $y=x^2$ tensor $\dd x\otimes\dd x$ ma zapis
$\frac1{4y}\dd y\otimes\dd y$, bo
$\dd x=\frac1{2\sqrt y}\dd y$. Zmiana współrzędnych nie
pozwala więc po prostu zastąpić symbolu $x$ symbolem $y$
przy pozostawieniu tych samych współczynników.
\end{przyklad}

Podobnie $\Lambda^k(T_pM)$ ma wymiar $\binom nk$.
W przypadku różniczki $\dd F_p\colon T_pM\to T_{F(p)}N$
między rozmaitościami tego samego wymiaru,
$\Lambda^n(\dd F_p)$ działa między ich jednowymiarowymi
potęgami zewnętrznymi. Jest niezerowe dokładnie wtedy,
gdy $\dd F_p$ jest izomorfizmem: wynika to z dowodu
Twierdzenia~\ref{thm:det-wlasnosci}, zastosowanego do baz
w dwóch różnych przestrzeniach. W mapach odpowiada temu
warunek niezerowego wyznacznika macierzy Jacobiego.
W odróżnieniu od endomorfizmu jednej przestrzeni, dla mapy
$T_pM\to T_{F(p)}N$ nie otrzymujemy jednak kanonicznego skalara
„wyznacznik” bez wyboru baz lub elementów objętości w obu włóknach.
Kanoniczne jest odwzorowanie $\Lambda^n(\dd F_p)$ między
dwiema prostymi; jego macierz liczbową określają wybrane mapy.

\begin{uwaga}[Cztery sposoby widzenia tej samej konstrukcji]
Iloczyn tensorowy powstaje z dwuliniowości
(Definicja~\ref{def:tensor-universal}); algebra tensorowa
gromadzi wszystkie stopnie i mnoży przez konkatenację
(Definicja~\ref{def:algebra-tensorowa}). Symetryzator
$\operatorname{Sym}_k$ wybiera część niezmienniczą względem
permutacji, zaś $\operatorname{Alt}_k$ część zmieniającą znak.
Algebra zewnętrzna jest ilorazem, w którym kwadrat każdego
wektora znika (Definicja~\ref{def:algebra-zewnetrzna}).
Twierdzenie~\ref{thm:wedge-alt} mówi dokładnie, jak przejść
od tego ilorazu do tensorów antysymetrycznych.
\end{uwaga}

\section{Różniczkowanie pól tensorowych: pochodna Liego}
\label{sec:lie-tensory}

Pole tensorowe ma wartości w różnych przestrzeniach nad
różnymi punktami. Dlatego wyrażenie ,,odejmij $T_q-T_p$''
nie ma sensu, dopóki nie wskażemy sposobu przeniesienia
obu tensorów do jednego włókna. W Sekcji~\ref{sec:liego}
przepływ pola $X$ zrobił to dla pól wektorowych;
ta sama konstrukcja działa na każdym typie tensora.

Samo różniczkowanie współczynników w mapie nie rozwiązuje
problemu: ich pochodne zmieniają się przy zmianie współrzędnych.
Na przykład na $(0,\infty)$ wektor $Y=\partial_x$ ma w
mapie $x$ stały współczynnik $Y^x=1$, więc
$\partial_xY^x=0$. Po podstawieniu $y=x^2$ mamy
$Y=2\sqrt y\,\partial_y$ i
$\partial_yY^y=1/\sqrt y\ne0$.
Różnica pochodzi od zmiany samej bazy $\partial_y$.
Pochodna Liego uwzględnia również tę zmianę.

\begin{definicja}[Cofnięcie pola tensorowego]
\index{cofniecie pola tensorowego@Cofnięcie pola tensorowego}
\label{def:cofniecie-tensora}
Niech $F\colon U\to V$ będzie dyfeomorfizmem między
otwartymi podzbiorami $M$. Dla pola wektorowego $Y$ i
pola kowektorowego $\alpha$ na $V$ definiujemy w $p\in U$
\begin{align*}
 (F^*Y)_p&:=(\dd F_p)^{-1}Y_{F(p)}\in T_pM,\\
 (F^*\alpha)_p(v)&:=\alpha_{F(p)}(\dd F_pv),
 \qquad v\in T_pM.
\end{align*}
Dla funkcji kładziemy $F^*f=f\circ F$.
Na tensorach prostych cofamy każdy czynnik osobno,
np.\ $F^*(Y\otimes\alpha)=F^*Y\otimes F^*\alpha$,
i przedłużamy liniowo na sumy w każdym włóknie. Do gładkości
wystarczy lokalne rozwinięcie w bazie tensorowej; nie zakładamy
istnienia jednej globalnej bazy ani globalnego rozkładu na tensory proste.
Dostajemy pole
tego samego typu na $U$. W szczególności
$F^*(fT)=(f\circ F)F^*T$, a nie $fF^*T$.
Współczynniki są gładkie: różniczka i jej macierz odwrotna
zależą gładko od punktu, a działania tensorowe są wieloliniowe.
\end{definicja}

Przy cofaniu wektora używamy odwrotnej różniczki,
a przy cofaniu kowektora różniczki bez odwracania:
obie wartości muszą ostatecznie należeć do włókna nad $p$.
Wzór dla tensorów prostych jest dobrze określony na
iloczynach tensorowych, ponieważ każda z operacji na
czynnikach jest liniowa. W szczególności cofnięcie zachowuje
iloczyn tensorowy i kontrakcję. To drugie można sprawdzić
na jednej parze:
\[
 (F^*\alpha)_p\bigl((F^*Y)_p\bigr)
 =\alpha_{F(p)}\bigl(\dd F_p(\dd F_p)^{-1}Y_{F(p)}\bigr)
 =\alpha_{F(p)}(Y_{F(p)}).
\]
Obie strony są cofnięciem funkcji $p\mapsto\alpha_p(Y_p)$;
z liniowości wynikają analogiczne równości dla innych tensorów.

\begin{definicja}[Pochodna Liego dowolnego pola tensorowego]
\index{pochodna liego dowolnego pola tensorowego@Pochodna Liego dowolnego pola tensorowego}
\label{def:lie-tensor}
Niech $\Phi_t=\Phi_t^X$ oznacza lokalny przepływ pola $X$.
Dla pola tensorowego $T$ dowolnego typu kładziemy
\begin{equation}\label{eq:lie-tensor-def}
 (\Lie_XT)_p
 :=\left.\frac{\dd}{\dd t}\right|_{t=0}
 (\Phi_t^*T)_p.
\end{equation}
Każdy tensor $(\Phi_t^*T)_p$ należy do tej samej
przestrzeni $T^r_s(T_pM)$, zatem pochodna jest
zwykłą pochodną krzywej w przestrzeni liniowej. Jest to gładkie
pole tensorowe: w każdej mapie współczynniki cofnięcia są gładkie
łącznie w $(t,p)$, więc ich pochodne po $t$ w zerze są gładkie w $p$.
\end{definicja}

W stopniu $(0,0)$ wzór~\eqref{eq:lie-tensor-def} daje
$\Lie_Xf=X(f)$, ponieważ $\Phi_t^*f=f\circ\Phi_t$.
Dla typu $(1,0)$ odzyskujemy dokładnie
$\Lie_XY=[X,Y]$ z Twierdzenia~\ref{thm:lie-nawias}.
Dowód tamtego twierdzenia wyjaśnia szczegółowo znak:
przenosimy $Y_{\Phi_t(p)}$ z powrotem różniczką
przepływu $\Phi_{-t}$.

\begin{twierdzenie}[Reguła Leibniza i zgodność z kontrakcją]
\label{thm:lie-tensor-reguly}
Dla pól tensorowych $S,T$, funkcji $f$, pola wektorowego
$Y$ i pola kowektorowego $\alpha$ mamy
\begin{align}
 \Lie_X(S\otimes T)
 &= (\Lie_XS)\otimes T+S\otimes(\Lie_XT),
 \label{eq:lie-tensor-leibniz}\\
 \Lie_X(fT)&=X(f)T+f\Lie_XT,
 \label{eq:lie-tensor-funkcja}\\
 (\Lie_X\alpha)(Y)
 &=X\bigl(\alpha(Y)\bigr)-\alpha([X,Y]),
 \label{eq:lie-kowektor}\\
 \Lie_X(C(S))&=C(\Lie_XS),
 \label{eq:lie-kontrakcja}
\end{align}
gdzie w ostatnim wzorze $C$ jest dowolną kontrakcją
jednego indeksu wektorowego z jednym kowektorowym.
\end{twierdzenie}
\begin{proof}
Ustalmy punkt $p$. Cofnięcie zachowuje iloczyn tensorowy:
$\Phi_t^*(S\otimes T)=\Phi_t^*S\otimes\Phi_t^*T$.
Obie krzywe tensorów po prawej stronie leżą w stałych
przestrzeniach nad $p$; różniczkując ich iloczyn
dwuliniowy w $t=0$ dostajemy
\eqref{eq:lie-tensor-leibniz}. Traktując $f$ jako tensor
typu $(0,0)$, otrzymujemy
\eqref{eq:lie-tensor-funkcja} z $\Lie_Xf=X(f)$.

Dla kowektora i wektora równanie o zachowaniu ewaluacji,
wykazane tuż przed definicją, przyjmuje postać
\[
 \Phi_t^*\bigl(\alpha(Y)\bigr)
 =(\Phi_t^*\alpha)(\Phi_t^*Y).
\]
Lewa strona po zróżniczkowaniu w zerze wynosi
$X(\alpha(Y))$. Prawa, przez zwykłą regułę
różniczkowania wartości kowektora na wektorze,
wynosi $(\Lie_X\alpha)(Y)+\alpha(\Lie_XY)$.
Z~$\Lie_XY=[X,Y]$ otrzymujemy
\eqref{eq:lie-kowektor}.

Wreszcie, cofnięcie i kontrakcja spełniają
$\Phi_t^*C(S)=C(\Phi_t^*S)$, co dla pojedynczej
pary sprawdziliśmy powyżej. Kontrakcja jest liniowym
odwzorowaniem między przestrzeniami tensorów w $p$, więc
można ją przenieść przez pochodną po $t$.
To daje~\eqref{eq:lie-kontrakcja}.
\end{proof}

\begin{twierdzenie}[Wzór we współrzędnych]
\label{thm:lie-tensor-wspolrzedne}
Niech $X=\sum_{a=1}^nX^a\partial_a$ i niech $T$ ma
współczynniki $T^I{}_J$ z~\eqref{eq:pole-tensorowe-lokalnie},
gdzie $I=(i_1,\ldots,i_r)$ i $J=(j_1,\ldots,j_s)$.
Symbol $I_{\mu\leftarrow a}$ oznacza zastąpienie w $I$
indeksu $i_\mu$ przez $a$, a $J_{\nu\leftarrow a}$
oznacza to samo w $J$. Wtedy
\begin{equation}\label{eq:lie-tensor-komponenty}
\begin{aligned}
 (\Lie_XT)^I{}_J
 &=\sum_{a=1}^nX^a\partial_a T^I{}_J
 -\sum_{\mu=1}^r\sum_{a=1}^n
   (\partial_aX^{i_\mu})T^{I_{\mu\leftarrow a}}{}_J\\
 &\quad+\sum_{\nu=1}^s\sum_{a=1}^n
   (\partial_{j_\nu}X^a)T^I{}_{J_{\nu\leftarrow a}}.
\end{aligned}
\end{equation}
Pierwszy wyraz różniczkuje funkcje współczynnikowe,
każdy indeks wektorowy wnosi wyraz z minusem,
a każdy indeks kowektorowy wyraz z plusem.
Na przykład dla pola operatorów $S\in\Gamma(T^1_1M)$
wzór ten ma łatwą do odczytania postać
\[
 (\Lie_XS)^i{}_j
 =\sum_a X^a\partial_aS^i{}_j
  -\sum_a S^a{}_j\partial_aX^i
  +\sum_a S^i{}_a\partial_jX^a.
\]
\end{twierdzenie}
\begin{proof}
Z~\eqref{eq:nawias-wspolrzedne}, zastosowanego do
$Y=\partial_i$, wynika
\[
 \Lie_X\partial_i=[X,\partial_i]
 =-\sum_a(\partial_iX^a)\partial_a,
\]
ponieważ współczynniki pola $\partial_i$ są stałe.
Z kolei $\dd x^j(\partial_i)=\delta_i^j$ jest funkcją
stałą. Wstawiając $Y=\partial_i$ i $\alpha=\dd x^j$
do~\eqref{eq:lie-kowektor}, dostajemy
\[
 (\Lie_X\dd x^j)(\partial_i)
 =X(\delta_i^j)-\dd x^j([X,\partial_i])
 =\partial_iX^j.
\]
Wektory $\partial_i$ tworzą bazę, zatem
$\Lie_X\dd x^j=\sum_i(\partial_iX^j)\dd x^i$.
Teraz rozwińmy $T$ według~\eqref{eq:pole-tensorowe-lokalnie}
i stosujmy~\eqref{eq:lie-tensor-leibniz} kolejno
do funkcji współczynnikowej oraz do każdego czynnika.
Różniczkowanie funkcji daje pierwszy wyraz
\eqref{eq:lie-tensor-komponenty}. Dla ustalonego miejsca
wektorowego $\mu$ współczynnik przy bazowym
$\partial_{i_\mu}$ po zróżniczkowaniu wcześniejszego
czynnika $\partial_a$ wynosi
$-\sum_a(\partial_aX^{i_\mu})T^{I_{\mu\leftarrow a}}{}_J$;
sumowanie po $\mu$ daje drugi wyraz wzoru.
Dla ustalonego miejsca kowektorowego $\nu$ wcześniejszy
czynnik $\dd x^a$ wnosi do współczynnika przy
$\dd x^{j_\nu}$ sumę
$\sum_a(\partial_{j_\nu}X^a)T^I{}_{J_{\nu\leftarrow a}}$;
po zsumowaniu po $\nu$ otrzymujemy trzeci wyraz.
Ponieważ tensory bazowe są liniowo niezależne,
ich współczynniki są dokładnie podanymi wyrażeniami.
\end{proof}

\begin{przyklad}[Sprawdzenie znaków na przepływie]
Na $\R^2$ weźmy $X=x\partial_x$. Jego przepływ to
$\Phi_t(x,y)=(e^tx,y)$, więc
\[
 \Phi_t^*\partial_x=e^{-t}\partial_x,
 \qquad \Phi_t^*\dd x=e^t\dd x.
\]
Po zróżniczkowaniu w zerze otrzymujemy
$\Lie_X\partial_x=-\partial_x=[X,\partial_x]$
oraz $\Lie_X\dd x=\dd x$. Pierwszy znak jest ujemny,
bo cofamy wektor różniczką odwrotnego przepływu;
drugi dodatni, bo kowektor działa na różniczce przepływu.
Wyniki są zgodne z tym, że $\dd x(\partial_x)=1$
pozostaje funkcją stałą: $1+(-1)=0$ w regule
Leibniza dla ewaluacji.
\end{przyklad}

\begin{przyklad}[Jak zmienia się tensor metryczny]
Dla $g=\dd x\otimes\dd x+\dd y\otimes\dd y$ i
$X=x\partial_x$ mamy
\[
 \Phi_t^*g=e^{2t}\dd x\otimes\dd x
             +\dd y\otimes\dd y,\qquad
 \Lie_Xg=2\dd x\otimes\dd x.
\]
Przepływ rozciąga kierunek $x$; gdy tensor mierzy
iloczyny skalarne dwóch wektorów, pojawia się czynnik
$e^t$ dla każdego z dwóch kowektorowych miejsc.
\end{przyklad}

\begin{uwaga}[Jaką pochodną zdefiniowaliśmy?]
Pochodna Liego porównuje tensory za pomocą przepływu
\emph{konkretnego} pola $X$. Jej wzór wymaga tylko struktury
gładkiej i pola $X$. Jeśli chcielibyśmy porównywać tensory
wzdłuż dowolnej krzywej bez wskazanego przepływu,
potrzebowalibyśmy dodatkowej reguły przenoszenia między
przestrzeniami stycznymi, nazywanej \emph{koneksją} (połączeniem).
Na przykład $\Lie_{x\partial_x}\partial_x=-\partial_x$
chociaż współczynniki pola $\partial_x$ są stałe:
wynik uwzględnia ruch bazy pod przepływem.
\end{uwaga}

\par\smallskip
{\footnotesize\noindent\emph{Dalsza lektura:}
\href{https://math.stanford.edu/~conrad/diffgeomPage/handouts/tensor.pdf}
{B. Conrad, \emph{Tensor algebras, tensor pairings, and duality}},
notatki do Math 396, s. 1--13: własności uniwersalne i dualność;
\href{https://doi.org/10.1007/978-1-4419-9982-5}
{J. M. Lee, \emph{Introduction to Smooth Manifolds}, wyd. 2},
rozdziały 12 i 14: pola tensorowe, formy alternujące
i pochodne Liego.\par}
\clearpage
